%% Ngân hàng đề Toán 12 (KNTT) -- Chương IV, Bài 11, Dạng 04
%% Dạng: Nguyên hàm thoả mãn điều kiện cho trước, bài toán thực tế
%% Phương án ĐẦU TIÊN trong \cauTN luôn là đáp án đúng (gói trondethi tự đảo).
%% Phần II: \y{nội dung}{d|s}. Phần III: đáp án đặt ở đối số cuối.

\cauTN[id=C4-B11-D04-NB-001,muc=NB,loigiai={$F(x)=4\sin x-3\mathrm{e}^x+C$. Vì $F(0)=5$ nên $C=8$. Vậy $F(x)=4\sin x-3\mathrm{e}^x+8$.}]{Tìm nguyên hàm $F(x)$ của hàm số $f(x)=4\cos x-3\mathrm{e}^x$ thoả mãn $F(0)=5$.}
  {$F(x)=4\sin x-3\mathrm{e}^x+8$}
  {$F(x)=4\sin x-3\mathrm{e}^x+5$}{$F(x)=4\sin x-3\mathrm{e}^x-8$}{$F(x)=4\sin x-3\mathrm{e}^x$}

\cauTN[id=C4-B11-D04-NB-002,muc=NB,loigiai={$F(x)=-\mathrm{e}^x-\dfrac{4x^{3}}{3}+C$. Vì $F(0)=5$ nên $C=6$. Vậy $F(x)=-\mathrm{e}^x-\dfrac{4x^{3}}{3}+6$.}]{Tìm nguyên hàm $F(x)$ của hàm số $f(x)=-\mathrm{e}^x-4x^{2}$ thoả mãn $F(0)=5$.}
  {$F(x)=-\mathrm{e}^x-\dfrac{4x^{3}}{3}+6$}
  {$F(x)=-\mathrm{e}^x-\dfrac{4x^{3}}{3}+5$}{$F(x)=-\mathrm{e}^x-\dfrac{4x^{3}}{3}-6$}{$F(x)=-\mathrm{e}^x-\dfrac{4x^{3}}{3}$}

\cauTN[id=C4-B11-D04-NB-003,muc=NB,loigiai={$F(x)=-\dfrac{3x^{2}}{2}-2x+C$. Vì $F(0)=-1$ nên $C=-1$. Vậy $F(x)=-\dfrac{3x^{2}}{2}-2x-1$.}]{Tìm nguyên hàm $F(x)$ của hàm số $f(x)=-3x-2$ thoả mãn $F(0)=-1$.}
  {$F(x)=-\dfrac{3x^{2}}{2}-2x-1$}
  {$F(x)=-\dfrac{3x^{2}}{2}-2x+1$}{$F(x)=-\dfrac{3x^{2}}{2}-2x$}{$F(x)=\dfrac{3x^{2}}{2}-2x-1$}

\cauTN[id=C4-B11-D04-NB-004,muc=NB,loigiai={$F(x)=\dfrac{x^{3}}{3}-4x^{2}+C$. Vì $F(0)=0$ nên $C=0$. Vậy $F(x)=\dfrac{x^{3}}{3}-4x^{2}$.}]{Tìm nguyên hàm $F(x)$ của hàm số $f(x)=x^{2}-8x$ thoả mãn $F(0)=0$.}
  {$F(x)=\dfrac{x^{3}}{3}-4x^{2}$}
  {$F(x)=\dfrac{x^{3}}{3}-4x^{2}+1$}{$F(x)=x^{3}-8x^{2}$}{$F(x)=2x-8$}

\cauTN[id=C4-B11-D04-NB-005,muc=NB,loigiai={$F(x)=-\dfrac{3x^{4}}{4}-4x^{2}+4x+C$. Vì $F(0)=4$ nên $C=4$. Vậy $F(x)=-\dfrac{3x^{4}}{4}-4x^{2}+4x+4$.}]{Tìm nguyên hàm $F(x)$ của hàm số $f(x)=-3x^{3}-8x+4$ thoả mãn $F(0)=4$.}
  {$F(x)=-\dfrac{3x^{4}}{4}-4x^{2}+4x+4$}
  {$F(x)=-\dfrac{3x^{4}}{4}-4x^{2}+4x-4$}{$F(x)=-\dfrac{3x^{4}}{4}-4x^{2}+4x$}{$F(x)=-\dfrac{3x^{4}}{4}-4x^{2}+4x+5$}

\cauTN[id=C4-B11-D04-NB-006,muc=NB,loigiai={$F(x)=\dfrac{3x^{2}}{2}+x+C$. Vì $F(0)=4$ nên $C=4$. Vậy $F(x)=\dfrac{3x^{2}}{2}+x+4$.}]{Tìm nguyên hàm $F(x)$ của hàm số $f(x)=3x+1$ thoả mãn $F(0)=4$.}
  {$F(x)=\dfrac{3x^{2}}{2}+x+4$}
  {$F(x)=\dfrac{3x^{2}}{2}+x-4$}{$F(x)=\dfrac{3x^{2}}{2}+x$}{$F(x)=\dfrac{3x^{2}}{2}+x+5$}

\cauTN[id=C4-B11-D04-NB-007,muc=NB,loigiai={$F(x)=-2\sin x-\dfrac{x^{2}}{2}+C$. Vì $F(0)=6$ nên $C=6$. Vậy $F(x)=-2\sin x-\dfrac{x^{2}}{2}+6$.}]{Tìm nguyên hàm $F(x)$ của hàm số $f(x)=-2\cos x-x$ thoả mãn $F(0)=6$.}
  {$F(x)=-2\sin x-\dfrac{x^{2}}{2}+6$}
  {$F(x)=-2\sin x-\dfrac{x^{2}}{2}-6$}{$F(x)=-2\sin x-\dfrac{x^{2}}{2}$}{$F(x)=-2\sin x-\dfrac{x^{2}}{2}+7$}

\cauTN[id=C4-B11-D04-NB-008,muc=NB,loigiai={$F(x)=2x^{2}+3x+C$. Vì $F(0)=-2$ nên $C=-2$. Vậy $F(x)=2x^{2}+3x-2$.}]{Tìm nguyên hàm $F(x)$ của hàm số $f(x)=4x+3$ thoả mãn $F(0)=-2$.}
  {$F(x)=2x^{2}+3x-2$}
  {$F(x)=2x^{2}+3x+2$}{$F(x)=2x^{2}+3x$}{$F(x)=2x^{2}+3x-1$}

\cauTN[id=C4-B11-D04-NB-009,muc=NB,loigiai={$F(x)=\dfrac{2x^{3}}{3}-2\sin x+C$. Vì $F(0)=1$ nên $C=1$. Vậy $F(x)=\dfrac{2x^{3}}{3}-2\sin x+1$.}]{Tìm nguyên hàm $F(x)$ của hàm số $f(x)=2x^{2}-2\cos x$ thoả mãn $F(0)=1$.}
  {$F(x)=\dfrac{2x^{3}}{3}-2\sin x+1$}
  {$F(x)=\dfrac{2x^{3}}{3}-2\sin x-1$}{$F(x)=\dfrac{2x^{3}}{3}-2\sin x$}{$F(x)=\dfrac{2x^{3}}{3}-2\sin x+2$}

\cauTN[id=C4-B11-D04-NB-010,muc=NB,loigiai={$F(x)=-\dfrac{4x^{3}}{3}-4x+C$. Vì $F(0)=6$ nên $C=6$. Vậy $F(x)=-\dfrac{4x^{3}}{3}-4x+6$.}]{Tìm nguyên hàm $F(x)$ của hàm số $f(x)=-4x^{2}-4$ thoả mãn $F(0)=6$.}
  {$F(x)=-\dfrac{4x^{3}}{3}-4x+6$}
  {$F(x)=-\dfrac{4x^{3}}{3}-4x-6$}{$F(x)=-\dfrac{4x^{3}}{3}-4x$}{$F(x)=-\dfrac{4x^{3}}{3}-4x+7$}

\cauTN[id=C4-B11-D04-NB-011,muc=NB,loigiai={$F(x)=-2x^{4}+2x^{2}+C$. Vì $F(0)=5$ nên $C=5$. Vậy $F(x)=-2x^{4}+2x^{2}+5$.}]{Tìm nguyên hàm $F(x)$ của hàm số $f(x)=-8x^{3}+4x$ thoả mãn $F(0)=5$.}
  {$F(x)=-2x^{4}+2x^{2}+5$}
  {$F(x)=-2x^{4}+2x^{2}-5$}{$F(x)=-2x^{4}+2x^{2}$}{$F(x)=-2x^{4}+2x^{2}+6$}

\cauTN[id=C4-B11-D04-NB-012,muc=NB,loigiai={$F(x)=-x^{3}+3\sin x+C$. Vì $F(0)=0$ nên $C=0$. Vậy $F(x)=-x^{3}+3\sin x$.}]{Tìm nguyên hàm $F(x)$ của hàm số $f(x)=-3x^{2}+3\cos x$ thoả mãn $F(0)=0$.}
  {$F(x)=-x^{3}+3\sin x$}
  {$F(x)=-x^{3}+3\sin x+1$}{$F(x)=-3x^{3}+3\sin x$}{$F(x)=-x^{3}+3\cos x$}

\cauTN[id=C4-B11-D04-NB-013,muc=NB,loigiai={$F(x)=x^{3}-\dfrac{3x^{2}}{2}+2x+C$. Vì $F(0)=-5$ nên $C=-5$. Vậy $F(x)=x^{3}-\dfrac{3x^{2}}{2}+2x-5$.}]{Tìm nguyên hàm $F(x)$ của hàm số $f(x)=3x^{2}-3x+2$ thoả mãn $F(0)=-5$.}
  {$F(x)=x^{3}-\dfrac{3x^{2}}{2}+2x-5$}
  {$F(x)=x^{3}-\dfrac{3x^{2}}{2}+2x+5$}{$F(x)=x^{3}-\dfrac{3x^{2}}{2}+2x$}{$F(x)=x^{3}-\dfrac{3x^{2}}{2}+2x-4$}

\cauTN[id=C4-B11-D04-NB-014,muc=NB,loigiai={$F(x)=2x^{3}-x+C$. Vì $F(0)=0$ nên $C=0$. Vậy $F(x)=2x^{3}-x$.}]{Tìm nguyên hàm $F(x)$ của hàm số $f(x)=6x^{2}-1$ thoả mãn $F(0)=0$.}
  {$F(x)=2x^{3}-x$}
  {$F(x)=2x^{3}-x+1$}{$F(x)=3x^{3}-x$}{$F(x)=12x$}

\cauTN[id=C4-B11-D04-NB-015,muc=NB,loigiai={$F(x)=3x^{4}-x^{3}+4x^{2}-2x+C$. Vì $F(0)=6$ nên $C=6$. Vậy $F(x)=3x^{4}-x^{3}+4x^{2}-2x+6$.}]{Tìm nguyên hàm $F(x)$ của hàm số $f(x)=12x^{3}-3x^{2}+8x-2$ thoả mãn $F(0)=6$.}
  {$F(x)=3x^{4}-x^{3}+4x^{2}-2x+6$}
  {$F(x)=3x^{4}-x^{3}+4x^{2}-2x-6$}{$F(x)=3x^{4}-x^{3}+4x^{2}-2x$}{$F(x)=3x^{4}-x^{3}+4x^{2}-2x+7$}

\cauTN[id=C4-B11-D04-NB-016,muc=NB,loigiai={$F(x)=4\sin x-\dfrac{x^{2}}{2}+C$. Vì $F(0)=0$ nên $C=0$. Vậy $F(x)=4\sin x-\dfrac{x^{2}}{2}$.}]{Tìm nguyên hàm $F(x)$ của hàm số $f(x)=4\cos x-x$ thoả mãn $F(0)=0$.}
  {$F(x)=4\sin x-\dfrac{x^{2}}{2}$}
  {$F(x)=4\sin x-\dfrac{x^{2}}{2}+1$}{$F(x)=4\cos x-\dfrac{x^{2}}{2}$}{$F(x)=-4\sin x-\dfrac{x^{2}}{2}$}

\cauTN[id=C4-B11-D04-NB-017,muc=NB,loigiai={$F(x)=-\dfrac{x^{4}}{2}+2x^{2}+C$. Vì $F(0)=0$ nên $C=0$. Vậy $F(x)=-\dfrac{x^{4}}{2}+2x^{2}$.}]{Tìm nguyên hàm $F(x)$ của hàm số $f(x)=-2x^{3}+4x$ thoả mãn $F(0)=0$.}
  {$F(x)=-\dfrac{x^{4}}{2}+2x^{2}$}
  {$F(x)=-\dfrac{x^{4}}{2}+2x^{2}+1$}{$F(x)=\dfrac{x^{4}}{2}+2x^{2}$}{$F(x)=-6x^{2}+4$}

\cauTN[id=C4-B11-D04-NB-018,muc=NB,loigiai={$F(x)=\dfrac{2x^{3}}{3}-2\sin x+C$. Vì $F(0)=4$ nên $C=4$. Vậy $F(x)=\dfrac{2x^{3}}{3}-2\sin x+4$.}]{Tìm nguyên hàm $F(x)$ của hàm số $f(x)=2x^{2}-2\cos x$ thoả mãn $F(0)=4$.}
  {$F(x)=\dfrac{2x^{3}}{3}-2\sin x+4$}
  {$F(x)=\dfrac{2x^{3}}{3}-2\sin x-4$}{$F(x)=\dfrac{2x^{3}}{3}-2\sin x$}{$F(x)=\dfrac{2x^{3}}{3}-2\sin x+5$}

\cauTN[id=C4-B11-D04-NB-019,muc=NB,loigiai={$F(x)=\sin x+3\cos x+C$. Vì $F(0)=-3$ nên $C=-6$. Vậy $F(x)=\sin x+3\cos x-6$.}]{Tìm nguyên hàm $F(x)$ của hàm số $f(x)=\cos x-3\sin x$ thoả mãn $F(0)=-3$.}
  {$F(x)=\sin x+3\cos x-6$}
  {$F(x)=\sin x+3\cos x-3$}{$F(x)=\sin x+3\cos x+6$}{$F(x)=\sin x+3\cos x$}

\cauTN[id=C4-B11-D04-NB-020,muc=NB,loigiai={$F(x)=-x^{3}+4x^{2}+x+C$. Vì $F(0)=-2$ nên $C=-2$. Vậy $F(x)=-x^{3}+4x^{2}+x-2$.}]{Tìm nguyên hàm $F(x)$ của hàm số $f(x)=-3x^{2}+8x+1$ thoả mãn $F(0)=-2$.}
  {$F(x)=-x^{3}+4x^{2}+x-2$}
  {$F(x)=-x^{3}+4x^{2}+x+2$}{$F(x)=-x^{3}+4x^{2}+x$}{$F(x)=-x^{3}+4x^{2}+x-1$}

\cauTN[id=C4-B11-D04-NB-021,muc=NB,loigiai={$F(x)=-x^{2}-2x+C$. Vì $F(0)=-1$ nên $C=-1$. Vậy $F(x)=-x^{2}-2x-1$.}]{Tìm nguyên hàm $F(x)$ của hàm số $f(x)=-2x-2$ thoả mãn $F(0)=-1$.}
  {$F(x)=-x^{2}-2x-1$}
  {$F(x)=-x^{2}-2x+1$}{$F(x)=-x^{2}-2x$}{$F(x)=-x^{2}+2x-1$}

\cauTN[id=C4-B11-D04-NB-022,muc=NB,loigiai={$F(x)=-3x^{4}+\dfrac{x^{3}}{3}+C$. Vì $F(0)=-5$ nên $C=-5$. Vậy $F(x)=-3x^{4}+\dfrac{x^{3}}{3}-5$.}]{Tìm nguyên hàm $F(x)$ của hàm số $f(x)=-12x^{3}+x^{2}$ thoả mãn $F(0)=-5$.}
  {$F(x)=-3x^{4}+\dfrac{x^{3}}{3}-5$}
  {$F(x)=-3x^{4}+\dfrac{x^{3}}{3}+5$}{$F(x)=-3x^{4}+\dfrac{x^{3}}{3}$}{$F(x)=-3x^{4}+\dfrac{x^{3}}{3}-4$}

\cauTN[id=C4-B11-D04-NB-023,muc=NB,loigiai={$F(x)=5\sin x-5\mathrm{e}^x+C$. Vì $F(0)=2$ nên $C=7$. Vậy $F(x)=5\sin x-5\mathrm{e}^x+7$.}]{Tìm nguyên hàm $F(x)$ của hàm số $f(x)=5\cos x-5\mathrm{e}^x$ thoả mãn $F(0)=2$.}
  {$F(x)=5\sin x-5\mathrm{e}^x+7$}
  {$F(x)=5\sin x-5\mathrm{e}^x+2$}{$F(x)=5\sin x-5\mathrm{e}^x-7$}{$F(x)=5\sin x-5\mathrm{e}^x$}

\cauTN[id=C4-B11-D04-NB-024,muc=NB,loigiai={$F(x)=-x^{3}+4x^{2}+3x+C$. Vì $F(0)=-2$ nên $C=-2$. Vậy $F(x)=-x^{3}+4x^{2}+3x-2$.}]{Tìm nguyên hàm $F(x)$ của hàm số $f(x)=-3x^{2}+8x+3$ thoả mãn $F(0)=-2$.}
  {$F(x)=-x^{3}+4x^{2}+3x-2$}
  {$F(x)=-x^{3}+4x^{2}+3x+2$}{$F(x)=-x^{3}+4x^{2}+3x$}{$F(x)=-x^{3}+4x^{2}+3x-1$}

\cauTN[id=C4-B11-D04-NB-025,muc=NB,loigiai={$F(x)=-\dfrac{x^{2}}{2}+\sin x+C$. Vì $F(0)=-1$ nên $C=-1$. Vậy $F(x)=-\dfrac{x^{2}}{2}+\sin x-1$.}]{Tìm nguyên hàm $F(x)$ của hàm số $f(x)=-x+\cos x$ thoả mãn $F(0)=-1$.}
  {$F(x)=-\dfrac{x^{2}}{2}+\sin x-1$}
  {$F(x)=-\dfrac{x^{2}}{2}+\sin x+1$}{$F(x)=-\dfrac{x^{2}}{2}+\sin x$}{$F(x)=-x^{2}+\sin x-1$}

\cauTN[id=C4-B11-D04-NB-026,muc=NB,loigiai={$F(x)=-3x^{3}-\dfrac{x^{2}}{2}-x+C$. Vì $F(0)=6$ nên $C=6$. Vậy $F(x)=-3x^{3}-\dfrac{x^{2}}{2}-x+6$.}]{Tìm nguyên hàm $F(x)$ của hàm số $f(x)=-9x^{2}-x-1$ thoả mãn $F(0)=6$.}
  {$F(x)=-3x^{3}-\dfrac{x^{2}}{2}-x+6$}
  {$F(x)=-3x^{3}-\dfrac{x^{2}}{2}-x-6$}{$F(x)=-3x^{3}-\dfrac{x^{2}}{2}-x$}{$F(x)=-3x^{3}-\dfrac{x^{2}}{2}-x+7$}

\cauTN[id=C4-B11-D04-NB-027,muc=NB,loigiai={$F(x)=-\dfrac{3x^{4}}{4}+x+C$. Vì $F(0)=0$ nên $C=0$. Vậy $F(x)=-\dfrac{3x^{4}}{4}+x$.}]{Tìm nguyên hàm $F(x)$ của hàm số $f(x)=-3x^{3}+1$ thoả mãn $F(0)=0$.}
  {$F(x)=-\dfrac{3x^{4}}{4}+x$}
  {$F(x)=-\dfrac{3x^{4}}{4}+x+1$}{$F(x)=-9x^{2}$}{$F(x)=-x^{4}+x$}

\cauTN[id=C4-B11-D04-NB-028,muc=NB,loigiai={$F(x)=-3x^{3}+\dfrac{x^{2}}{2}+C$. Vì $F(0)=3$ nên $C=3$. Vậy $F(x)=-3x^{3}+\dfrac{x^{2}}{2}+3$.}]{Tìm nguyên hàm $F(x)$ của hàm số $f(x)=-9x^{2}+x$ thoả mãn $F(0)=3$.}
  {$F(x)=-3x^{3}+\dfrac{x^{2}}{2}+3$}
  {$F(x)=-3x^{3}+\dfrac{x^{2}}{2}-3$}{$F(x)=-3x^{3}+\dfrac{x^{2}}{2}$}{$F(x)=-3x^{3}+\dfrac{x^{2}}{2}+4$}

\cauTN[id=C4-B11-D04-NB-029,muc=NB,loigiai={$F(x)=4x+\dfrac{3x^{2}}{2}+C$. Vì $F(0)=5$ nên $C=5$. Vậy $F(x)=4x+\dfrac{3x^{2}}{2}+5$.}]{Tìm nguyên hàm $F(x)$ của hàm số $f(x)=4+3x$ thoả mãn $F(0)=5$.}
  {$F(x)=4x+\dfrac{3x^{2}}{2}+5$}
  {$F(x)=4x+\dfrac{3x^{2}}{2}-5$}{$F(x)=4x+\dfrac{3x^{2}}{2}$}{$F(x)=4x+\dfrac{3x^{2}}{2}+6$}

\cauTN[id=C4-B11-D04-NB-030,muc=NB,loigiai={$F(x)=x^{3}+x^{2}+3x+C$. Vì $F(0)=-1$ nên $C=-1$. Vậy $F(x)=x^{3}+x^{2}+3x-1$.}]{Tìm nguyên hàm $F(x)$ của hàm số $f(x)=3x^{2}+2x+3$ thoả mãn $F(0)=-1$.}
  {$F(x)=x^{3}+x^{2}+3x-1$}
  {$F(x)=x^{3}+x^{2}+3x+1$}{$F(x)=x^{3}+x^{2}+3x$}{$F(x)=-x^{3}+x^{2}+3x-1$}

\cauTN[id=C4-B11-D04-NB-031,muc=NB,loigiai={$F(x)=\dfrac{4x^{3}}{3}+5\sin x+C$. Vì $F(0)=5$ nên $C=5$. Vậy $F(x)=\dfrac{4x^{3}}{3}+5\sin x+5$.}]{Tìm nguyên hàm $F(x)$ của hàm số $f(x)=4x^{2}+5\cos x$ thoả mãn $F(0)=5$.}
  {$F(x)=\dfrac{4x^{3}}{3}+5\sin x+5$}
  {$F(x)=\dfrac{4x^{3}}{3}+5\sin x-5$}{$F(x)=\dfrac{4x^{3}}{3}+5\sin x$}{$F(x)=\dfrac{4x^{3}}{3}+5\sin x+6$}

\cauTN[id=C4-B11-D04-NB-032,muc=NB,loigiai={$F(x)=-x^{4}+4x^{2}+x+C$. Vì $F(0)=0$ nên $C=0$. Vậy $F(x)=-x^{4}+4x^{2}+x$.}]{Tìm nguyên hàm $F(x)$ của hàm số $f(x)=-4x^{3}+8x+1$ thoả mãn $F(0)=0$.}
  {$F(x)=-x^{4}+4x^{2}+x$}
  {$F(x)=-x^{4}+4x^{2}+x+1$}{$F(x)=-x^{4}+4x^{2}-x$}{$F(x)=-4x^{4}+8x^{2}+x$}

\cauTN[id=C4-B11-D04-NB-033,muc=NB,loigiai={$F(x)=-3x^{2}-4x+C$. Vì $F(0)=3$ nên $C=3$. Vậy $F(x)=-3x^{2}-4x+3$.}]{Tìm nguyên hàm $F(x)$ của hàm số $f(x)=-6x-4$ thoả mãn $F(0)=3$.}
  {$F(x)=-3x^{2}-4x+3$}
  {$F(x)=-3x^{2}-4x-3$}{$F(x)=-3x^{2}-4x$}{$F(x)=-3x^{2}-4x+4$}

\cauTN[id=C4-B11-D04-NB-034,muc=NB,loigiai={$F(x)=3x^{3}-2x^{2}-4x+C$. Vì $F(0)=6$ nên $C=6$. Vậy $F(x)=3x^{3}-2x^{2}-4x+6$.}]{Tìm nguyên hàm $F(x)$ của hàm số $f(x)=9x^{2}-4x-4$ thoả mãn $F(0)=6$.}
  {$F(x)=3x^{3}-2x^{2}-4x+6$}
  {$F(x)=3x^{3}-2x^{2}-4x-6$}{$F(x)=3x^{3}-2x^{2}-4x$}{$F(x)=3x^{3}-2x^{2}-4x+7$}

\cauTN[id=C4-B11-D04-NB-035,muc=NB,loigiai={$F(x)=-4x^{3}+2x^{2}-x+C$. Vì $F(0)=0$ nên $C=0$. Vậy $F(x)=-4x^{3}+2x^{2}-x$.}]{Tìm nguyên hàm $F(x)$ của hàm số $f(x)=-12x^{2}+4x-1$ thoả mãn $F(0)=0$.}
  {$F(x)=-4x^{3}+2x^{2}-x$}
  {$F(x)=-4x^{3}+2x^{2}-x+1$}{$F(x)=4x^{3}+2x^{2}-x$}{$F(x)=-4x^{3}-2x^{2}-x$}

\cauTN[id=C4-B11-D04-NB-036,muc=NB,loigiai={$F(x)=-2x^{2}+3x+C$. Vì $F(0)=3$ nên $C=3$. Vậy $F(x)=-2x^{2}+3x+3$.}]{Tìm nguyên hàm $F(x)$ của hàm số $f(x)=-4x+3$ thoả mãn $F(0)=3$.}
  {$F(x)=-2x^{2}+3x+3$}
  {$F(x)=-2x^{2}+3x-3$}{$F(x)=-2x^{2}+3x$}{$F(x)=-2x^{2}+3x+4$}

\cauDS[id=C4-B11-D04-NB-037,muc=NB,loigiai={$F(x)=-\dfrac{x^{4}}{2}-x^{3}+\dfrac{3x^{2}}{2}+4x+C$, $F(0)=1\Rightarrow C=1$, nên $F(x)=-\dfrac{x^{4}}{2}-x^{3}+\dfrac{3x^{2}}{2}+4x+1$. a) Đúng. b) Đúng. c) Sai: giá trị đúng là $-1$. d) Đúng.}]{Cho hàm số $f(x)=-2x^{3}-3x^{2}+3x+4$ và $F(x)$ là nguyên hàm của $f(x)$ thoả mãn $F(0)=1$. Xét tính đúng sai của các mệnh đề sau:}{%
  \y{$F(2)=-1$.}{d}
  \y{$\displaystyle\int f(x)\,\mathrm{d}x=-\dfrac{x^{4}}{2}-x^{3}+\dfrac{3x^{2}}{2}+4x+C$.}{d}
  \y{$F(2)=2$.}{s}
  \y{$F(x)=-\dfrac{x^{4}}{2}-x^{3}+\dfrac{3x^{2}}{2}+4x+1$.}{d}}

\cauDS[id=C4-B11-D04-NB-038,muc=NB,loigiai={$F(x)=2x^{4}+4x^{2}+C$, $F(0)=6\Rightarrow C=6$, nên $F(x)=2x^{4}+4x^{2}+6$. a) Đúng. b) Sai: giá trị đúng là $12$. c) Sai. d) Sai.}]{Cho hàm số $f(x)=8x^{3}+8x$ và $F(x)$ là nguyên hàm của $f(x)$ thoả mãn $F(0)=6$. Xét tính đúng sai của các mệnh đề sau:}{%
  \y{$F'(x)=8x^{3}+8x$.}{d}
  \y{$F(-1)=10$.}{s}
  \y{$F(x)=2x^{4}+4x^{2}$.}{s}
  \y{$\displaystyle\int f(x)\,\mathrm{d}x=2x^{4}-4x^{2}+C$.}{s}}

\cauDS[id=C4-B11-D04-NB-039,muc=NB,loigiai={$F(x)=2x^{2}-x+C$, $F(0)=-2\Rightarrow C=-2$, nên $F(x)=2x^{2}-x-2$. a) Đúng. b) Sai: giá trị đúng là $4$. c) Đúng. d) Đúng.}]{Cho hàm số $f(x)=4x-1$ và $F(x)$ là nguyên hàm của $f(x)$ thoả mãn $F(0)=-2$. Xét tính đúng sai của các mệnh đề sau:}{%
  \y{$\displaystyle\int f(x)\,\mathrm{d}x=2x^{2}-x+C$.}{d}
  \y{$F(2)=7$.}{s}
  \y{$F(2)=4$.}{d}
  \y{$F'(x)=4x-1$.}{d}}

\cauDS[id=C4-B11-D04-NB-040,muc=NB,loigiai={$F(x)=-\dfrac{x^{3}}{3}-\dfrac{x^{2}}{2}+2x+C$, $F(0)=-3\Rightarrow C=-3$, nên $F(x)=-\dfrac{x^{3}}{3}-\dfrac{x^{2}}{2}+2x-3$. a) Đúng. b) Đúng. c) Sai. d) Sai.}]{Cho hàm số $f(x)=-x^{2}-x+2$ và $F(x)$ là nguyên hàm của $f(x)$ thoả mãn $F(0)=-3$. Xét tính đúng sai của các mệnh đề sau:}{%
  \y{$F(1)=-\dfrac{11}{6}$.}{d}
  \y{$F'(x)=-x^{2}-x+2$.}{d}
  \y{$F(x)=-\dfrac{x^{3}}{3}-\dfrac{x^{2}}{2}+2x$.}{s}
  \y{$\displaystyle\int f(x)\,\mathrm{d}x=-\dfrac{x^{3}}{3}+\dfrac{x^{2}}{2}+2x+C$.}{s}}

\cauDS[id=C4-B11-D04-NB-041,muc=NB,loigiai={$F(x)=-4x^{2}-3x+C$, $F(0)=-3\Rightarrow C=-3$, nên $F(x)=-4x^{2}-3x-3$. a) Sai: giá trị đúng là $-25$. b) Đúng. c) Sai. d) Đúng.}]{Cho hàm số $f(x)=-8x-3$ và $F(x)$ là nguyên hàm của $f(x)$ thoả mãn $F(0)=-3$. Xét tính đúng sai của các mệnh đề sau:}{%
  \y{$F(2)=-24$.}{s}
  \y{$F(2)=-25$.}{d}
  \y{$\displaystyle\int f(x)\,\mathrm{d}x=-16x^{2}-3x+C$.}{s}
  \y{$F'(x)=-8x-3$.}{d}}

\cauDS[id=C4-B11-D04-NB-042,muc=NB,loigiai={$F(x)=-3x^{2}-4x+C$, $F(0)=0\Rightarrow C=0$, nên $F(x)=-3x^{2}-4x$. a) Sai: giá trị đúng là $1$. b) Đúng. c) Sai. d) Đúng.}]{Cho hàm số $f(x)=-6x-4$ và $F(x)$ là nguyên hàm của $f(x)$ thoả mãn $F(0)=0$. Xét tính đúng sai của các mệnh đề sau:}{%
  \y{$F(-1)=4$.}{s}
  \y{$\displaystyle\int f(x)\,\mathrm{d}x=-3x^{2}-4x+C$.}{d}
  \y{$\displaystyle\int f(x)\,\mathrm{d}x=-3x^{2}-4+C$.}{s}
  \y{$F(x)=-3x^{2}-4x$.}{d}}

\cauDS[id=C4-B11-D04-NB-043,muc=NB,loigiai={$F(x)=3x^{2}+4x+C$, $F(0)=-5\Rightarrow C=-5$, nên $F(x)=3x^{2}+4x-5$. a) Đúng. b) Đúng. c) Sai. d) Sai: giá trị đúng là $2$.}]{Cho hàm số $f(x)=6x+4$ và $F(x)$ là nguyên hàm của $f(x)$ thoả mãn $F(0)=-5$. Xét tính đúng sai của các mệnh đề sau:}{%
  \y{$F(1)=2$.}{d}
  \y{$\displaystyle\int f(x)\,\mathrm{d}x=3x^{2}+4x+C$.}{d}
  \y{$F(x)=3x^{2}+4x$.}{s}
  \y{$F(1)=3$.}{s}}

\cauDS[id=C4-B11-D04-NB-044,muc=NB,loigiai={$F(x)=-\dfrac{3x^{2}}{2}+3x+C$, $F(0)=-5\Rightarrow C=-5$, nên $F(x)=-\dfrac{3x^{2}}{2}+3x-5$. a) Đúng. b) Đúng. c) Sai: giá trị đúng là $-\dfrac{19}{2}$. d) Sai.}]{Cho hàm số $f(x)=-3x+3$ và $F(x)$ là nguyên hàm của $f(x)$ thoả mãn $F(0)=-5$. Xét tính đúng sai của các mệnh đề sau:}{%
  \y{$F(x)=-\dfrac{3x^{2}}{2}+3x-5$.}{d}
  \y{$\displaystyle\int f(x)\,\mathrm{d}x=-\dfrac{3x^{2}}{2}+3x+C$.}{d}
  \y{$F(-1)=-\dfrac{17}{2}$.}{s}
  \y{$\displaystyle\int f(x)\,\mathrm{d}x=\dfrac{3x^{2}}{2}+3x+C$.}{s}}

\cauDS[id=C4-B11-D04-NB-045,muc=NB,loigiai={$F(x)=-2x^{3}+x^{2}-x+C$, $F(0)=4\Rightarrow C=4$, nên $F(x)=-2x^{3}+x^{2}-x+4$. a) Sai: giá trị đúng là $2$. b) Đúng. c) Đúng. d) Sai.}]{Cho hàm số $f(x)=-6x^{2}+2x-1$ và $F(x)$ là nguyên hàm của $f(x)$ thoả mãn $F(0)=4$. Xét tính đúng sai của các mệnh đề sau:}{%
  \y{$F(1)=0$.}{s}
  \y{$\displaystyle\int f(x)\,\mathrm{d}x=-2x^{3}+x^{2}-x+C$.}{d}
  \y{$F'(x)=-6x^{2}+2x-1$.}{d}
  \y{$\displaystyle\int f(x)\,\mathrm{d}x=-2x^{3}+x^{2}+x+C$.}{s}}

\cauDS[id=C4-B11-D04-NB-046,muc=NB,loigiai={$F(x)=4x^{4}-2x^{3}-4x^{2}+2x+C$, $F(0)=3\Rightarrow C=3$, nên $F(x)=4x^{4}-2x^{3}-4x^{2}+2x+3$. a) Sai: giá trị đúng là $3$. b) Đúng. c) Đúng. d) Sai.}]{Cho hàm số $f(x)=16x^{3}-6x^{2}-8x+2$ và $F(x)$ là nguyên hàm của $f(x)$ thoả mãn $F(0)=3$. Xét tính đúng sai của các mệnh đề sau:}{%
  \y{$F(-1)=1$.}{s}
  \y{$F'(x)=16x^{3}-6x^{2}-8x+2$.}{d}
  \y{$\displaystyle\int f(x)\,\mathrm{d}x=4x^{4}-2x^{3}-4x^{2}+2x+C$.}{d}
  \y{$\displaystyle\int f(x)\,\mathrm{d}x=4x^{4}-2x^{3}-4x^{2}+2+C$.}{s}}

\cauDS[id=C4-B11-D04-NB-047,muc=NB,loigiai={$F(x)=-2x^{4}-x+C$, $F(0)=0\Rightarrow C=0$, nên $F(x)=-2x^{4}-x$. a) Đúng. b) Đúng. c) Sai: giá trị đúng là $-3$. d) Sai.}]{Cho hàm số $f(x)=-8x^{3}-1$ và $F(x)$ là nguyên hàm của $f(x)$ thoả mãn $F(0)=0$. Xét tính đúng sai của các mệnh đề sau:}{%
  \y{$F(1)=-3$.}{d}
  \y{$\displaystyle\int f(x)\,\mathrm{d}x=-2x^{4}-x+C$.}{d}
  \y{$F(1)=0$.}{s}
  \y{$\displaystyle\int f(x)\,\mathrm{d}x=-8x^{4}-x+C$.}{s}}

\cauDS[id=C4-B11-D04-NB-048,muc=NB,loigiai={$F(x)=-3x^{4}+\dfrac{x^{3}}{3}-x^{2}+C$, $F(0)=-1\Rightarrow C=-1$, nên $F(x)=-3x^{4}+\dfrac{x^{3}}{3}-x^{2}-1$. a) Sai. b) Sai: giá trị đúng là $-\dfrac{151}{3}$. c) Đúng. d) Đúng.}]{Cho hàm số $f(x)=-12x^{3}+x^{2}-2x$ và $F(x)$ là nguyên hàm của $f(x)$ thoả mãn $F(0)=-1$. Xét tính đúng sai của các mệnh đề sau:}{%
  \y{$\displaystyle\int f(x)\,\mathrm{d}x=-12x^{4}+x^{3}-2x^{2}+C$.}{s}
  \y{$F(2)=-\dfrac{142}{3}$.}{s}
  \y{$F'(x)=-12x^{3}+x^{2}-2x$.}{d}
  \y{$\displaystyle\int f(x)\,\mathrm{d}x=-3x^{4}+\dfrac{x^{3}}{3}-x^{2}+C$.}{d}}

\cauTLN[id=C4-B11-D04-NB-049,muc=NB,loigiai={$F(x)=-3x^{2}-4x-3$ nên $F(-1)=-2$.}]{Biết $F(x)$ là một nguyên hàm của hàm số $f(x)=-6x-4$ và $F(0)=-3$. Tính $F(-1)$.}{-2}

\cauTLN[id=C4-B11-D04-NB-050,muc=NB,loigiai={$F(x)=-x^{4}-x^{2}-1$ nên $F(-1)=-3$.}]{Biết $F(x)$ là một nguyên hàm của hàm số $f(x)=-4x^{3}-2x$ và $F(0)=-1$. Tính $F(-1)$.}{-3}

\cauTLN[id=C4-B11-D04-NB-051,muc=NB,loigiai={$F(x)=-2x^{4}-2x^{3}+3x$ nên $F(1)=-1$.}]{Biết $F(x)$ là một nguyên hàm của hàm số $f(x)=-8x^{3}-6x^{2}+3$ và $F(0)=0$. Tính $F(1)$.}{-1}

\cauTLN[id=C4-B11-D04-NB-052,muc=NB,loigiai={$F(x)=2x^{3}+3x^{2}+x-1$ nên $F(2)=29$.}]{Biết $F(x)$ là một nguyên hàm của hàm số $f(x)=6x^{2}+6x+1$ và $F(0)=-1$. Tính $F(2)$.}{29}

\cauTLN[id=C4-B11-D04-NB-053,muc=NB,loigiai={$F(x)=x^{3}+\dfrac{3x^{2}}{2}+3x-4$ nên $F(3)=\dfrac{91}{2}$.}]{Biết $F(x)$ là một nguyên hàm của hàm số $f(x)=3x^{2}+3x+3$ và $F(0)=-4$. Tính $F(3)$.}{45,5}

\cauTLN[id=C4-B11-D04-NB-054,muc=NB,loigiai={$F(x)=x^{3}-2x^{2}-4x-5$ nên $F(2)=-13$.}]{Biết $F(x)$ là một nguyên hàm của hàm số $f(x)=3x^{2}-4x-4$ và $F(0)=-5$. Tính $F(2)$.}{-13}

\cauTLN[id=C4-B11-D04-NB-055,muc=NB,loigiai={$F(x)=2x^{3}-2x^{2}$ nên $F(3)=36$.}]{Biết $F(x)$ là một nguyên hàm của hàm số $f(x)=6x^{2}-4x$ và $F(0)=0$. Tính $F(3)$.}{36}

\cauTLN[id=C4-B11-D04-NB-056,muc=NB,loigiai={$F(x)=x^{3}-\dfrac{x^{2}}{2}-3x-5$ nên $F(2)=-5$.}]{Biết $F(x)$ là một nguyên hàm của hàm số $f(x)=3x^{2}-x-3$ và $F(0)=-5$. Tính $F(2)$.}{-5}

\cauTLN[id=C4-B11-D04-NB-057,muc=NB,loigiai={$F(x)=-3x^{4}+4x^{3}-x^{2}+x+4$ nên $F(-1)=-5$.}]{Biết $F(x)$ là một nguyên hàm của hàm số $f(x)=-12x^{3}+12x^{2}-2x+1$ và $F(0)=4$. Tính $F(-1)$.}{-5}

\cauTLN[id=C4-B11-D04-NB-058,muc=NB,loigiai={$F(x)=-4x^{4}+x^{2}+2x-3$ nên $F(-1)=-8$.}]{Biết $F(x)$ là một nguyên hàm của hàm số $f(x)=-16x^{3}+2x+2$ và $F(0)=-3$. Tính $F(-1)$.}{-8}

\cauTLN[id=C4-B11-D04-NB-059,muc=NB,loigiai={$F(x)=x^{3}-x^{2}+2x-5$ nên $F(1)=-3$.}]{Biết $F(x)$ là một nguyên hàm của hàm số $f(x)=3x^{2}-2x+2$ và $F(0)=-5$. Tính $F(1)$.}{-3}

\cauTLN[id=C4-B11-D04-NB-060,muc=NB,loigiai={$F(x)=-2x^{4}-\dfrac{4x^{3}}{3}-3$ nên $F(3)=-201$.}]{Biết $F(x)$ là một nguyên hàm của hàm số $f(x)=-8x^{3}-4x^{2}$ và $F(0)=-3$. Tính $F(3)$.}{-201}

\cauTN[id=C4-B11-D04-TH-001,muc=TH,loigiai={$F(x)=5\sin x+\dfrac{x^{2}}{2}+C$; $F(0)=4\Rightarrow C=4$. Vậy $F(1)=5 \sin{(1 )} + \dfrac{9}{2}$.}]{Biết $F(x)$ là một nguyên hàm của hàm số $f(x)=5\cos x+x$ và $F(0)=4$. Giá trị của $F(1)$ bằng}
  {$5 \sin{(1 )} + \dfrac{9}{2}$}
  {$\dfrac{1}{2} + 5 \sin{(1 )}$}{$- \dfrac{7}{2} + 5 \sin{(1 )}$}{$5 \sin{(1 )} + \dfrac{11}{2}$}

\cauTN[id=C4-B11-D04-TH-002,muc=TH,loigiai={Lượng nước $V(t)=\displaystyle\int r(t)\,\mathrm{d}t=3t^2+30t$ (vì $V(0)=0$); $V(5)=225$ (lít).}]{Nước được bơm vào một bể với tốc độ $r(t)=6t+30$ (lít/phút), $t$ tính bằng phút. Lúc đầu bể không có nước. Sau $5$ phút, lượng nước trong bể là}
  {$225$ lít}
  {$60$ lít}{$300$ lít}{$75$ lít}

\cauTN[id=C4-B11-D04-TH-003,muc=TH,loigiai={$v(t)=\displaystyle\int a(t)\,\mathrm{d}t=\dfrac{3}{2}t^2+4t+C$, $v(0)=2\Rightarrow C=2$. $v(3)=\dfrac{55}{2}$ (m/s).}]{Một vật chuyển động với gia tốc $a(t)=3t+4$ (m/s$^2$), $t$ tính bằng giây. Biết vận tốc ban đầu của vật là $2$ m/s. Vận tốc của vật tại thời điểm $t=3$ giây là}
  {$27{,}5$ m/s}
  {$15$ m/s}{$41$ m/s}{$25{,}5$ m/s}

\cauTN[id=C4-B11-D04-TH-004,muc=TH,loigiai={$F(x)=-5x-\dfrac{5}{x}+C$; $F(1)=-4\Rightarrow C=6$. Vậy $F(2)=-\dfrac{13}{2}$.}]{Biết $F(x)$ là một nguyên hàm của hàm số $f(x)=-5+\dfrac{5}{x^2}$ trên khoảng $(0;+\infty)$ và $F(1)=-4$. Giá trị của $F(2)$ bằng}
  {$-\dfrac{13}{2}$}
  {$-\dfrac{25}{2}$}{$-\dfrac{37}{2}$}{$-\dfrac{11}{2}$}

\cauTN[id=C4-B11-D04-TH-005,muc=TH,loigiai={Lượng nước $V(t)=\displaystyle\int r(t)\,\mathrm{d}t=2t^2+10t$ (vì $V(0)=0$); $V(15)=600$ (lít).}]{Nước được bơm vào một bể với tốc độ $r(t)=4t+10$ (lít/phút), $t$ tính bằng phút. Lúc đầu bể không có nước. Sau $15$ phút, lượng nước trong bể là}
  {$600$ lít}
  {$70$ lít}{$1050$ lít}{$450$ lít}

\cauTN[id=C4-B11-D04-TH-006,muc=TH,loigiai={Gọi $S(t)$ là quãng đường, $S'(t)=v(t)$ nên $S(t)=t^{2} + 6 t+C$; $S(0)=0\Rightarrow C=0$. $S(30)=1080$ (m).}]{Một máy bay chạy đà trên đường băng với vận tốc $v(t)=6+2t$ (m/s), trong đó $t$ là thời gian (giây) kể từ lúc bắt đầu chạy đà. Sau $30$ giây máy bay rời đường băng. Quãng đường (mét) máy bay đã chạy trên đường băng là}
  {$1080$ m}
  {$1980$ m}{$66$ m}{$900$ m}

\cauTN[id=C4-B11-D04-TH-007,muc=TH,loigiai={Gọi $S(t)$ là quãng đường, $S'(t)=v(t)$ nên $S(t)=t^{2} + 5 t+C$; $S(0)=0\Rightarrow C=0$. $S(30)=1050$ (m).}]{Một máy bay chạy đà trên đường băng với vận tốc $v(t)=5+2t$ (m/s), trong đó $t$ là thời gian (giây) kể từ lúc bắt đầu chạy đà. Sau $30$ giây máy bay rời đường băng. Quãng đường (mét) máy bay đã chạy trên đường băng là}
  {$1050$ m}
  {$1950$ m}{$65$ m}{$900$ m}

\cauTN[id=C4-B11-D04-TH-008,muc=TH,loigiai={$F(x)=-4\sin x-4\mathrm{e}^x+C$; $F(0)=-5\Rightarrow C=-1$. Vậy $F(1)=- 4 \mathrm{e} - 4 \sin{(1 )} - 1$.}]{Biết $F(x)$ là một nguyên hàm của hàm số $f(x)=-4\cos x-4\mathrm{e}^x$ và $F(0)=-5$. Giá trị của $F(1)$ bằng}
  {$- 4 \mathrm{e} - 4 \sin{(1 )} - 1$}
  {$- 4 \mathrm{e} - 4 \sin{(1 )}$}{$- 4 \mathrm{e} - 4 \sin{(1 )} + 1$}{$- 4 \mathrm{e} - 4 \sin{(1 )} + 4$}

\cauTN[id=C4-B11-D04-TH-009,muc=TH,loigiai={Gọi $S(t)$ là quãng đường, $S'(t)=v(t)$ nên $S(t)=t^{2} + 5 t+C$; $S(0)=0\Rightarrow C=0$. $S(40)=1800$ (m).}]{Một máy bay chạy đà trên đường băng với vận tốc $v(t)=5+2t$ (m/s), trong đó $t$ là thời gian (giây) kể từ lúc bắt đầu chạy đà. Sau $40$ giây máy bay rời đường băng. Quãng đường (mét) máy bay đã chạy trên đường băng là}
  {$1800$ m}
  {$3400$ m}{$85$ m}{$1600$ m}

\cauTN[id=C4-B11-D04-TH-010,muc=TH,loigiai={Gọi $S(t)$ là quãng đường, $S'(t)=v(t)$ nên $S(t)=2 t^{2} + 4 t+C$; $S(0)=0\Rightarrow C=0$. $S(20)=880$ (m).}]{Một máy bay chạy đà trên đường băng với vận tốc $v(t)=4+4t$ (m/s), trong đó $t$ là thời gian (giây) kể từ lúc bắt đầu chạy đà. Sau $20$ giây máy bay rời đường băng. Quãng đường (mét) máy bay đã chạy trên đường băng là}
  {$880$ m}
  {$1680$ m}{$84$ m}{$800$ m}

\cauTN[id=C4-B11-D04-TH-011,muc=TH,loigiai={$R(x)=200x-\dfrac{1}{20}x^2$; $R(100)=19500$ (triệu đồng).}]{Tốc độ biến động doanh thu khi bán $x$ sản phẩm là $R'(x)=200-\dfrac{1}{10}x$ (triệu đồng/sản phẩm). Biết $R(0)=0$. Doanh thu khi bán $100$ sản phẩm là}
  {$19500$ triệu đồng}
  {$19000$ triệu đồng}{$190$ triệu đồng}{$20000$ triệu đồng}

\cauTN[id=C4-B11-D04-TH-012,muc=TH,loigiai={$R(x)=300x-\dfrac{1}{40}x^2$; $R(1000)=275000$ (triệu đồng).}]{Tốc độ biến động doanh thu khi bán $x$ sản phẩm là $R'(x)=300-\dfrac{1}{20}x$ (triệu đồng/sản phẩm). Biết $R(0)=0$. Doanh thu khi bán $1000$ sản phẩm là}
  {$275000$ triệu đồng}
  {$250000$ triệu đồng}{$250$ triệu đồng}{$300000$ triệu đồng}

\cauTN[id=C4-B11-D04-TH-013,muc=TH,loigiai={$F(x)=-5x-3\cos x+C$; $F(0)=5\Rightarrow C=8$. Vậy $F(2)=-2 - 3 \cos{(2 )}$.}]{Biết $F(x)$ là một nguyên hàm của hàm số $f(x)=-5+3\sin x$ và $F(0)=5$. Giá trị của $F(2)$ bằng}
  {$-2 - 3 \cos{(2 )}$}
  {$-10 - 3 \cos{(2 )}$}{$-18 - 3 \cos{(2 )}$}{$-1 - 3 \cos{(2 )}$}

\cauTN[id=C4-B11-D04-TH-014,muc=TH,loigiai={Gọi $S(t)$ là quãng đường, $S'(t)=v(t)$ nên $S(t)=t^{2} + 4 t+C$; $S(0)=0\Rightarrow C=0$. $S(40)=1760$ (m).}]{Một máy bay chạy đà trên đường băng với vận tốc $v(t)=4+2t$ (m/s), trong đó $t$ là thời gian (giây) kể từ lúc bắt đầu chạy đà. Sau $40$ giây máy bay rời đường băng. Quãng đường (mét) máy bay đã chạy trên đường băng là}
  {$1760$ m}
  {$3360$ m}{$84$ m}{$1600$ m}

\cauTN[id=C4-B11-D04-TH-015,muc=TH,loigiai={Gọi $S(t)$ là quãng đường, $S'(t)=v(t)$ nên $S(t)=2 t^{2} + 8 t+C$; $S(0)=0\Rightarrow C=0$. $S(40)=3520$ (m).}]{Một máy bay chạy đà trên đường băng với vận tốc $v(t)=8+4t$ (m/s), trong đó $t$ là thời gian (giây) kể từ lúc bắt đầu chạy đà. Sau $40$ giây máy bay rời đường băng. Quãng đường (mét) máy bay đã chạy trên đường băng là}
  {$3520$ m}
  {$6720$ m}{$168$ m}{$3200$ m}

\cauTN[id=C4-B11-D04-TH-016,muc=TH,loigiai={$F(x)=-2x^{2}+\dfrac{2x^{3}}{3}+C$; $F(0)=-3\Rightarrow C=-3$. Vậy $F(1)=-\dfrac{13}{3}$.}]{Biết $F(x)$ là một nguyên hàm của hàm số $f(x)=-4x+2x^{2}$ và $F(0)=-3$. Giá trị của $F(1)$ bằng}
  {$-\dfrac{13}{3}$}
  {$-\dfrac{4}{3}$}{$\dfrac{5}{3}$}{$-\dfrac{10}{3}$}

\cauTN[id=C4-B11-D04-TH-017,muc=TH,loigiai={$P(t)=500+20t+\dfrac{3}{2}t^2$ (vì $P(0)=500$). $P(12)=956$ (con).}]{Một quần thể vi khuẩn ban đầu có $500$ con. Tốc độ tăng trưởng của quần thể tại thời điểm $t$ (giờ) là $P'(t)=20+3t$ (con/giờ). Số vi khuẩn sau $12$ giờ là}
  {$956$ con}
  {$1172$ con}{$456$ con}{$556$ con}

\cauTN[id=C4-B11-D04-TH-018,muc=TH,loigiai={$P(t)=1000+50t+2t^2$ (vì $P(0)=1000$). $P(10)=1700$ (con).}]{Một quần thể vi khuẩn ban đầu có $1000$ con. Tốc độ tăng trưởng của quần thể tại thời điểm $t$ (giờ) là $P'(t)=50+4t$ (con/giờ). Số vi khuẩn sau $10$ giờ là}
  {$1700$ con}
  {$1900$ con}{$700$ con}{$1090$ con}

\cauTN[id=C4-B11-D04-TH-019,muc=TH,loigiai={$v(t)=\displaystyle\int a(t)\,\mathrm{d}t=\dfrac{3}{2}t^2+2t+C$, $v(0)=10\Rightarrow C=10$. $v(2)=20$ (m/s).}]{Một vật chuyển động với gia tốc $a(t)=3t+2$ (m/s$^2$), $t$ tính bằng giây. Biết vận tốc ban đầu của vật là $10$ m/s. Vận tốc của vật tại thời điểm $t=2$ giây là}
  {$20$ m/s}
  {$18$ m/s}{$26$ m/s}{$10$ m/s}

\cauTN[id=C4-B11-D04-TH-020,muc=TH,loigiai={$F(x)=\dfrac{1}{x}+2\mathrm{e}^x+C$; $F(1)=5\Rightarrow C=4 - 2 \mathrm{e}$. Vậy $F(2)=- 2 \mathrm{e} + \dfrac{9}{2} + 2 \mathrm{e}^{2}$.}]{Biết $F(x)$ là một nguyên hàm của hàm số $f(x)=-\dfrac{1}{x^2}+2\mathrm{e}^x$ trên khoảng $(0;+\infty)$ và $F(1)=5$. Giá trị của $F(2)$ bằng}
  {$- 2 \mathrm{e} + \dfrac{9}{2} + 2 \mathrm{e}^{2}$}
  {$\dfrac{1}{2} + 2 \mathrm{e}^{2}$}{$- \dfrac{7}{2} + 2 \mathrm{e} + 2 \mathrm{e}^{2}$}{$- 2 \mathrm{e} + \dfrac{11}{2} + 2 \mathrm{e}^{2}$}

\cauTN[id=C4-B11-D04-TH-021,muc=TH,loigiai={$F(x)=-\dfrac{3}{x}-3\cos x+C$; $F(1)=5\Rightarrow C=3 \cos{(1 )} + 8$. Vậy $F(2)=- 3 \cos{(2 )} + 3 \cos{(1 )} + \dfrac{13}{2}$.}]{Biết $F(x)$ là một nguyên hàm của hàm số $f(x)=\dfrac{3}{x^2}+3\sin x$ trên khoảng $(0;+\infty)$ và $F(1)=5$. Giá trị của $F(2)$ bằng}
  {$- 3 \cos{(2 )} + 3 \cos{(1 )} + \dfrac{13}{2}$}
  {$- \dfrac{3}{2} - 3 \cos{(2 )}$}{$- \dfrac{19}{2} - 3 \cos{(1 )} - 3 \cos{(2 )}$}{$- 3 \cos{(2 )} + 3 \cos{(1 )} + \dfrac{15}{2}$}

\cauTN[id=C4-B11-D04-TH-022,muc=TH,loigiai={Lượng nước $V(t)=\displaystyle\int r(t)\,\mathrm{d}t=3t^2+20t$ (vì $V(0)=0$); $V(15)=975$ (lít).}]{Nước được bơm vào một bể với tốc độ $r(t)=6t+20$ (lít/phút), $t$ tính bằng phút. Lúc đầu bể không có nước. Sau $15$ phút, lượng nước trong bể là}
  {$975$ lít}
  {$110$ lít}{$1650$ lít}{$675$ lít}

\cauTN[id=C4-B11-D04-TH-023,muc=TH,loigiai={$F(x)=\dfrac{3x^{2}}{2}+\dfrac{2x^{3}}{3}+C$; $F(0)=5\Rightarrow C=5$. Vậy $F(1)=\dfrac{43}{6}$.}]{Biết $F(x)$ là một nguyên hàm của hàm số $f(x)=3x+2x^{2}$ và $F(0)=5$. Giá trị của $F(1)$ bằng}
  {$\dfrac{43}{6}$}
  {$\dfrac{13}{6}$}{$-\dfrac{17}{6}$}{$\dfrac{49}{6}$}

\cauTN[id=C4-B11-D04-TH-024,muc=TH,loigiai={Gọi $S(t)$ là quãng đường, $S'(t)=v(t)$ nên $S(t)=2 t^{2} + 5 t+C$; $S(0)=0\Rightarrow C=0$. $S(25)=1375$ (m).}]{Một máy bay chạy đà trên đường băng với vận tốc $v(t)=5+4t$ (m/s), trong đó $t$ là thời gian (giây) kể từ lúc bắt đầu chạy đà. Sau $25$ giây máy bay rời đường băng. Quãng đường (mét) máy bay đã chạy trên đường băng là}
  {$1375$ m}
  {$2625$ m}{$105$ m}{$1250$ m}

\cauTN[id=C4-B11-D04-TH-025,muc=TH,loigiai={$F(x)=\dfrac{2x^{3}}{3}-3\sin x+C$; $F(0)=3\Rightarrow C=3$. Vậy $F(2)=\dfrac{25}{3} - 3 \sin{(2 )}$.}]{Biết $F(x)$ là một nguyên hàm của hàm số $f(x)=2x^{2}-3\cos x$ và $F(0)=3$. Giá trị của $F(2)$ bằng}
  {$\dfrac{25}{3} - 3 \sin{(2 )}$}
  {$\dfrac{16}{3} - 3 \sin{(2 )}$}{$\dfrac{7}{3} - 3 \sin{(2 )}$}{$\dfrac{28}{3} - 3 \sin{(2 )}$}

\cauTN[id=C4-B11-D04-TH-026,muc=TH,loigiai={$F(x)=x^{3}-\mathrm{e}^x+C$; $F(0)=-2\Rightarrow C=-1$. Vậy $F(1)=- \mathrm{e}$.}]{Biết $F(x)$ là một nguyên hàm của hàm số $f(x)=3x^{2}-\mathrm{e}^x$ và $F(0)=-2$. Giá trị của $F(1)$ bằng}
  {$- \mathrm{e}$}
  {$1 - \mathrm{e}$}{$2 - \mathrm{e}$}{$- \mathrm{e} - 2$}

\cauTN[id=C4-B11-D04-TH-027,muc=TH,loigiai={$F(x)=3\mathrm{e}^x-\dfrac{2}{x}+C$; $F(1)=4\Rightarrow C=6 - 3 \mathrm{e}$. Vậy $F(2)=- 3 \mathrm{e} + 5 + 3 \mathrm{e}^{2}$.}]{Biết $F(x)$ là một nguyên hàm của hàm số $f(x)=3\mathrm{e}^x+\dfrac{2}{x^2}$ trên khoảng $(0;+\infty)$ và $F(1)=4$. Giá trị của $F(2)$ bằng}
  {$- 3 \mathrm{e} + 5 + 3 \mathrm{e}^{2}$}
  {$-1 + 3 \mathrm{e}^{2}$}{$-7 + 3 \mathrm{e} + 3 \mathrm{e}^{2}$}{$- 3 \mathrm{e} + 6 + 3 \mathrm{e}^{2}$}

\cauTN[id=C4-B11-D04-TH-028,muc=TH,loigiai={$F(x)=-2x-\dfrac{4}{x}+C$; $F(1)=-4\Rightarrow C=2$. Vậy $F(2)=-4$.}]{Biết $F(x)$ là một nguyên hàm của hàm số $f(x)=-2+\dfrac{4}{x^2}$ trên khoảng $(0;+\infty)$ và $F(1)=-4$. Giá trị của $F(2)$ bằng}
  {$-4$}
  {$-6$}{$-8$}{$-3$}

\cauTN[id=C4-B11-D04-TH-029,muc=TH,loigiai={$v(t)=\displaystyle\int a(t)\,\mathrm{d}t=3t^2+2t+C$, $v(0)=2\Rightarrow C=2$. $v(4)=58$ (m/s).}]{Một vật chuyển động với gia tốc $a(t)=6t+2$ (m/s$^2$), $t$ tính bằng giây. Biết vận tốc ban đầu của vật là $2$ m/s. Vận tốc của vật tại thời điểm $t=4$ giây là}
  {$58$ m/s}
  {$28$ m/s}{$106$ m/s}{$56$ m/s}

\cauTN[id=C4-B11-D04-TH-030,muc=TH,loigiai={$F(x)=2x+\dfrac{3x^{2}}{2}+C$; $F(0)=1\Rightarrow C=1$. Vậy $F(2)=11$.}]{Biết $F(x)$ là một nguyên hàm của hàm số $f(x)=2+3x$ và $F(0)=1$. Giá trị của $F(2)$ bằng}
  {$11$}
  {$10$}{$9$}{$12$}

\cauTN[id=C4-B11-D04-TH-031,muc=TH,loigiai={$F(x)=3\mathrm{e}^x+\dfrac{5x^{2}}{2}+C$; $F(0)=-2\Rightarrow C=-5$. Vậy $F(2)=5 + 3 \mathrm{e}^{2}$.}]{Biết $F(x)$ là một nguyên hàm của hàm số $f(x)=3\mathrm{e}^x+5x$ và $F(0)=-2$. Giá trị của $F(2)$ bằng}
  {$5 + 3 \mathrm{e}^{2}$}
  {$10 + 3 \mathrm{e}^{2}$}{$15 + 3 \mathrm{e}^{2}$}{$6 + 3 \mathrm{e}^{2}$}

\cauTN[id=C4-B11-D04-TH-032,muc=TH,loigiai={$P(t)=2000+40t+2t^2$ (vì $P(0)=2000$). $P(12)=2768$ (con).}]{Một quần thể vi khuẩn ban đầu có $2000$ con. Tốc độ tăng trưởng của quần thể tại thời điểm $t$ (giờ) là $P'(t)=40+4t$ (con/giờ). Số vi khuẩn sau $12$ giờ là}
  {$2768$ con}
  {$3056$ con}{$768$ con}{$2088$ con}

\cauTN[id=C4-B11-D04-TH-033,muc=TH,loigiai={$F(x)=4\sin x+5\cos x+C$; $F(0)=0\Rightarrow C=-5$. Vậy $F(1)=-5 + 5 \cos{(1 )} + 4 \sin{(1 )}$.}]{Biết $F(x)$ là một nguyên hàm của hàm số $f(x)=4\cos x-5\sin x$ và $F(0)=0$. Giá trị của $F(1)$ bằng}
  {$-5 + 5 \cos{(1 )} + 4 \sin{(1 )}$}
  {$5 \cos{(1 )} + 4 \sin{(1 )}$}{$5 \cos{(1 )} + 4 \sin{(1 )} + 5$}{$-4 + 5 \cos{(1 )} + 4 \sin{(1 )}$}

\cauTN[id=C4-B11-D04-TH-034,muc=TH,loigiai={Gọi $S(t)$ là quãng đường, $S'(t)=v(t)$ nên $S(t)=\dfrac{3 t^{2}}{2} + 8 t+C$; $S(0)=0\Rightarrow C=0$. $S(40)=2720$ (m).}]{Một máy bay chạy đà trên đường băng với vận tốc $v(t)=8+3t$ (m/s), trong đó $t$ là thời gian (giây) kể từ lúc bắt đầu chạy đà. Sau $40$ giây máy bay rời đường băng. Quãng đường (mét) máy bay đã chạy trên đường băng là}
  {$2720$ m}
  {$5120$ m}{$128$ m}{$2400$ m}

\cauTN[id=C4-B11-D04-TH-035,muc=TH,loigiai={Xe dừng khi $v(t)=0\Leftrightarrow t=4$. $S(t)=\displaystyle\int v(t)\,\mathrm{d}t=-\dfrac{5}{2}t^2+20t$ (vì $S(0)=0$), $S(4)=40$ (m).}]{Một ô tô đang chạy với vận tốc $20$ m/s thì người lái đạp phanh; từ thời điểm đó ô tô chuyển động chậm dần đều với vận tốc $v(t)=-5t+20$ (m/s), $t$ là thời gian (giây) tính từ lúc đạp phanh. Từ lúc đạp phanh đến khi dừng hẳn, ô tô di chuyển được bao nhiêu mét?}
  {$40$ m}
  {$80$ m}{$4$ m}{$60$ m}

\cauTN[id=C4-B11-D04-TH-036,muc=TH,loigiai={Xe dừng khi $v(t)=0\Leftrightarrow t=3$. $S(t)=\displaystyle\int v(t)\,\mathrm{d}t=-2t^2+12t$ (vì $S(0)=0$), $S(3)=18$ (m).}]{Một ô tô đang chạy với vận tốc $12$ m/s thì người lái đạp phanh; từ thời điểm đó ô tô chuyển động chậm dần đều với vận tốc $v(t)=-4t+12$ (m/s), $t$ là thời gian (giây) tính từ lúc đạp phanh. Từ lúc đạp phanh đến khi dừng hẳn, ô tô di chuyển được bao nhiêu mét?}
  {$18$ m}
  {$36$ m}{$3$ m}{$30$ m}

\cauTN[id=C4-B11-D04-TH-037,muc=TH,loigiai={Lượng nước $V(t)=\displaystyle\int r(t)\,\mathrm{d}t=3t^2+10t$ (vì $V(0)=0$); $V(15)=825$ (lít).}]{Nước được bơm vào một bể với tốc độ $r(t)=6t+10$ (lít/phút), $t$ tính bằng phút. Lúc đầu bể không có nước. Sau $15$ phút, lượng nước trong bể là}
  {$825$ lít}
  {$100$ lít}{$1500$ lít}{$675$ lít}

\cauTN[id=C4-B11-D04-TH-038,muc=TH,loigiai={Lượng nước $V(t)=\displaystyle\int r(t)\,\mathrm{d}t=2t^2+10t$ (vì $V(0)=0$); $V(10)=300$ (lít).}]{Nước được bơm vào một bể với tốc độ $r(t)=4t+10$ (lít/phút), $t$ tính bằng phút. Lúc đầu bể không có nước. Sau $10$ phút, lượng nước trong bể là}
  {$300$ lít}
  {$50$ lít}{$500$ lít}{$200$ lít}

\cauTN[id=C4-B11-D04-TH-039,muc=TH,loigiai={$F(x)=2x^{2}-4\sin x+C$; $F(0)=-5\Rightarrow C=-5$. Vậy $F(1)=- 4 \sin{(1 )} - 3$.}]{Biết $F(x)$ là một nguyên hàm của hàm số $f(x)=4x-4\cos x$ và $F(0)=-5$. Giá trị của $F(1)$ bằng}
  {$- 4 \sin{(1 )} - 3$}
  {$2 - 4 \sin{(1 )}$}{$7 - 4 \sin{(1 )}$}{$- 4 \sin{(1 )} - 2$}

\cauTN[id=C4-B11-D04-TH-040,muc=TH,loigiai={$F(x)=\cos x+5\sin x+C$; $F(0)=-2\Rightarrow C=-3$. Vậy $F(1)=-3 + \cos{(1 )} + 5 \sin{(1 )}$.}]{Biết $F(x)$ là một nguyên hàm của hàm số $f(x)=-\sin x+5\cos x$ và $F(0)=-2$. Giá trị của $F(1)$ bằng}
  {$-3 + \cos{(1 )} + 5 \sin{(1 )}$}
  {$\cos{(1 )} + 5 \sin{(1 )}$}{$\cos{(1 )} + 3 + 5 \sin{(1 )}$}{$-2 + \cos{(1 )} + 5 \sin{(1 )}$}

\cauTN[id=C4-B11-D04-TH-041,muc=TH,loigiai={$R(x)=200x-\dfrac{1}{20}x^2$; $R(1000)=150000$ (triệu đồng).}]{Tốc độ biến động doanh thu khi bán $x$ sản phẩm là $R'(x)=200-\dfrac{1}{10}x$ (triệu đồng/sản phẩm). Biết $R(0)=0$. Doanh thu khi bán $1000$ sản phẩm là}
  {$150000$ triệu đồng}
  {$100000$ triệu đồng}{$100$ triệu đồng}{$200000$ triệu đồng}

\cauTN[id=C4-B11-D04-TH-042,muc=TH,loigiai={$R(x)=300x-\dfrac{1}{40}x^2$; $R(500)=143750$ (triệu đồng).}]{Tốc độ biến động doanh thu khi bán $x$ sản phẩm là $R'(x)=300-\dfrac{1}{20}x$ (triệu đồng/sản phẩm). Biết $R(0)=0$. Doanh thu khi bán $500$ sản phẩm là}
  {$143750$ triệu đồng}
  {$137500$ triệu đồng}{$275$ triệu đồng}{$150000$ triệu đồng}

\cauTN[id=C4-B11-D04-TH-043,muc=TH,loigiai={Gọi $S(t)$ là quãng đường, $S'(t)=v(t)$ nên $S(t)=\dfrac{3 t^{2}}{2} + 8 t+C$; $S(0)=0\Rightarrow C=0$. $S(30)=1590$ (m).}]{Một máy bay chạy đà trên đường băng với vận tốc $v(t)=8+3t$ (m/s), trong đó $t$ là thời gian (giây) kể từ lúc bắt đầu chạy đà. Sau $30$ giây máy bay rời đường băng. Quãng đường (mét) máy bay đã chạy trên đường băng là}
  {$1590$ m}
  {$2940$ m}{$98$ m}{$1350$ m}

\cauTN[id=C4-B11-D04-TH-044,muc=TH,loigiai={$P(t)=2000+30t+2t^2$ (vì $P(0)=2000$). $P(10)=2500$ (con).}]{Một quần thể vi khuẩn ban đầu có $2000$ con. Tốc độ tăng trưởng của quần thể tại thời điểm $t$ (giờ) là $P'(t)=30+4t$ (con/giờ). Số vi khuẩn sau $10$ giờ là}
  {$2500$ con}
  {$2700$ con}{$500$ con}{$2070$ con}

\cauTN[id=C4-B11-D04-TH-045,muc=TH,loigiai={$F(x)=\sin x-x^{2}+C$; $F(0)=3\Rightarrow C=3$. Vậy $F(2)=-1 + \sin{(2 )}$.}]{Biết $F(x)$ là một nguyên hàm của hàm số $f(x)=\cos x-2x$ và $F(0)=3$. Giá trị của $F(2)$ bằng}
  {$-1 + \sin{(2 )}$}
  {$-4 + \sin{(2 )}$}{$-7 + \sin{(2 )}$}{$\sin{(2 )}$}

\cauTN[id=C4-B11-D04-TH-046,muc=TH,loigiai={$R(x)=300x-\dfrac{1}{20}x^2$; $R(500)=137500$ (triệu đồng).}]{Tốc độ biến động doanh thu khi bán $x$ sản phẩm là $R'(x)=300-\dfrac{1}{10}x$ (triệu đồng/sản phẩm). Biết $R(0)=0$. Doanh thu khi bán $500$ sản phẩm là}
  {$137500$ triệu đồng}
  {$125000$ triệu đồng}{$250$ triệu đồng}{$150000$ triệu đồng}

\cauTN[id=C4-B11-D04-TH-047,muc=TH,loigiai={$F(x)=-\dfrac{4}{x}+5x+C$; $F(1)=-5\Rightarrow C=-6$. Vậy $F(4)=13$.}]{Biết $F(x)$ là một nguyên hàm của hàm số $f(x)=\dfrac{4}{x^2}+5$ trên khoảng $(0;+\infty)$ và $F(1)=-5$. Giá trị của $F(4)$ bằng}
  {$13$}
  {$19$}{$25$}{$14$}

\cauTN[id=C4-B11-D04-TH-048,muc=TH,loigiai={$F(x)=-\dfrac{x^{2}}{2}+2\mathrm{e}^x+C$; $F(0)=-5\Rightarrow C=-7$. Vậy $F(2)=-9 + 2 \mathrm{e}^{2}$.}]{Biết $F(x)$ là một nguyên hàm của hàm số $f(x)=-x+2\mathrm{e}^x$ và $F(0)=-5$. Giá trị của $F(2)$ bằng}
  {$-9 + 2 \mathrm{e}^{2}$}
  {$-2 + 2 \mathrm{e}^{2}$}{$5 + 2 \mathrm{e}^{2}$}{$-8 + 2 \mathrm{e}^{2}$}

\cauDS[id=C4-B11-D04-TH-049,muc=TH,loigiai={$s(t)=\displaystyle\int v(t)\,\mathrm{d}t=t^{3}+t^{2}$ (vì $s(0)=0$); $a(t)=v'(t)=6t+2$. a) Sai: đó là gia tốc. b) Đúng. c) Sai: đó là vận tốc tại $t=1$. d) Sai: giá trị đúng là $80$ m.}]{Một vật chuyển động thẳng với vận tốc $v(t)=3t^{2}+2t$ (m/s), $t$ là thời gian tính bằng giây, quãng đường tính từ thời điểm $t=0$. Xét tính đúng sai của các mệnh đề sau:}{%
  \y{Quãng đường vật đi được sau $t$ giây là $s(t)=6t+2$ (m).}{s}
  \y{Vận tốc trung bình của vật trong $4$ giây đầu là $20$ m/s.}{d}
  \y{Gia tốc của vật tại thời điểm $t=1$ giây là $5$ m/s$^2$.}{s}
  \y{Quãng đường vật đi được trong $4$ giây đầu là $84$ m.}{s}}

\cauDS[id=C4-B11-D04-TH-050,muc=TH,loigiai={$s(t)=\displaystyle\int v(t)\,\mathrm{d}t=2t^{3}+3t^{2}$ (vì $s(0)=0$); $a(t)=v'(t)=12t+6$. a) Sai: đó là gia tốc. b) Sai: đó là vận tốc tại $t=1$. c) Sai: giá trị đúng là $81$ m. d) Đúng.}]{Một vật chuyển động thẳng với vận tốc $v(t)=6t^{2}+6t$ (m/s), $t$ là thời gian tính bằng giây, quãng đường tính từ thời điểm $t=0$. Xét tính đúng sai của các mệnh đề sau:}{%
  \y{Quãng đường vật đi được sau $t$ giây là $s(t)=12t+6$ (m).}{s}
  \y{Gia tốc của vật tại thời điểm $t=1$ giây là $12$ m/s$^2$.}{s}
  \y{Quãng đường vật đi được trong $3$ giây đầu là $84$ m.}{s}
  \y{Gia tốc của vật tại thời điểm $t=1$ giây là $18$ m/s$^2$.}{d}}

\cauDS[id=C4-B11-D04-TH-051,muc=TH,loigiai={$s(t)=\displaystyle\int v(t)\,\mathrm{d}t=2t^{3}+3t^{2}$ (vì $s(0)=0$); $a(t)=v'(t)=12t+6$. a) Đúng. b) Sai: giá trị đúng là $176$ m. c) Sai: đó là vận tốc tại $t=1$. d) Đúng.}]{Một vật chuyển động thẳng với vận tốc $v(t)=6t^{2}+6t$ (m/s), $t$ là thời gian tính bằng giây, quãng đường tính từ thời điểm $t=0$. Xét tính đúng sai của các mệnh đề sau:}{%
  \y{Quãng đường vật đi được trong $4$ giây đầu là $176$ m.}{d}
  \y{Quãng đường vật đi được trong $4$ giây đầu là $180$ m.}{s}
  \y{Gia tốc của vật tại thời điểm $t=1$ giây là $12$ m/s$^2$.}{s}
  \y{Vận tốc trung bình của vật trong $4$ giây đầu là $44$ m/s.}{d}}

\cauDS[id=C4-B11-D04-TH-052,muc=TH,loigiai={$s(t)=\displaystyle\int v(t)\,\mathrm{d}t=t^{3}+2t^{2}$ (vì $s(0)=0$); $a(t)=v'(t)=6t+4$. a) Sai: giá trị đúng là $96$ m. b) Đúng. c) Sai: đó là gia tốc. d) Đúng.}]{Một vật chuyển động thẳng với vận tốc $v(t)=3t^{2}+4t$ (m/s), $t$ là thời gian tính bằng giây, quãng đường tính từ thời điểm $t=0$. Xét tính đúng sai của các mệnh đề sau:}{%
  \y{Quãng đường vật đi được trong $4$ giây đầu là $100$ m.}{s}
  \y{Vận tốc trung bình của vật trong $4$ giây đầu là $24$ m/s.}{d}
  \y{Quãng đường vật đi được sau $t$ giây là $s(t)=6t+4$ (m).}{s}
  \y{Gia tốc của vật tại thời điểm $t=1$ giây là $10$ m/s$^2$.}{d}}

\cauDS[id=C4-B11-D04-TH-053,muc=TH,loigiai={$s(t)=\displaystyle\int v(t)\,\mathrm{d}t=t^{3}+3t^{2}$ (vì $s(0)=0$); $a(t)=v'(t)=6t+6$. a) Sai: đó là gia tốc. b) Sai: giá trị đúng là $112$ m. c) Sai: đó là vận tốc tại $t=1$. d) Đúng.}]{Một vật chuyển động thẳng với vận tốc $v(t)=3t^{2}+6t$ (m/s), $t$ là thời gian tính bằng giây, quãng đường tính từ thời điểm $t=0$. Xét tính đúng sai của các mệnh đề sau:}{%
  \y{Quãng đường vật đi được sau $t$ giây là $s(t)=6t+6$ (m).}{s}
  \y{Quãng đường vật đi được trong $4$ giây đầu là $116$ m.}{s}
  \y{Gia tốc của vật tại thời điểm $t=1$ giây là $9$ m/s$^2$.}{s}
  \y{Gia tốc của vật tại thời điểm $t=1$ giây là $12$ m/s$^2$.}{d}}

\cauDS[id=C4-B11-D04-TH-054,muc=TH,loigiai={$s(t)=\displaystyle\int v(t)\,\mathrm{d}t=t^{3}+2t^{2}$ (vì $s(0)=0$); $a(t)=v'(t)=6t+4$. a) Đúng. b) Sai: đó là gia tốc. c) Đúng. d) Đúng.}]{Một vật chuyển động thẳng với vận tốc $v(t)=3t^{2}+4t$ (m/s), $t$ là thời gian tính bằng giây, quãng đường tính từ thời điểm $t=0$. Xét tính đúng sai của các mệnh đề sau:}{%
  \y{Gia tốc của vật tại thời điểm $t=1$ giây là $10$ m/s$^2$.}{d}
  \y{Quãng đường vật đi được sau $t$ giây là $s(t)=6t+4$ (m).}{s}
  \y{Quãng đường vật đi được sau $t$ giây là $s(t)=t^{3}+2t^{2}$ (m).}{d}
  \y{Vận tốc trung bình của vật trong $4$ giây đầu là $24$ m/s.}{d}}

\cauDS[id=C4-B11-D04-TH-055,muc=TH,loigiai={$s(t)=\displaystyle\int v(t)\,\mathrm{d}t=t^{3}+2t^{2}$ (vì $s(0)=0$); $a(t)=v'(t)=6t+4$. a) Sai: đó là vận tốc tại $t=1$. b) Đúng. c) Đúng. d) Sai: đó là gia tốc.}]{Một vật chuyển động thẳng với vận tốc $v(t)=3t^{2}+4t$ (m/s), $t$ là thời gian tính bằng giây, quãng đường tính từ thời điểm $t=0$. Xét tính đúng sai của các mệnh đề sau:}{%
  \y{Gia tốc của vật tại thời điểm $t=1$ giây là $7$ m/s$^2$.}{s}
  \y{Gia tốc của vật tại thời điểm $t=1$ giây là $10$ m/s$^2$.}{d}
  \y{Quãng đường vật đi được sau $t$ giây là $s(t)=t^{3}+2t^{2}$ (m).}{d}
  \y{Quãng đường vật đi được sau $t$ giây là $s(t)=6t+4$ (m).}{s}}

\cauDS[id=C4-B11-D04-TH-056,muc=TH,loigiai={$s(t)=\displaystyle\int v(t)\,\mathrm{d}t=2t^{3}+2t^{2}$ (vì $s(0)=0$); $a(t)=v'(t)=12t+4$. a) Sai: giá trị đúng là $160$ m. b) Sai: đó là gia tốc. c) Sai: đó là vận tốc tại $t=1$. d) Đúng.}]{Một vật chuyển động thẳng với vận tốc $v(t)=6t^{2}+4t$ (m/s), $t$ là thời gian tính bằng giây, quãng đường tính từ thời điểm $t=0$. Xét tính đúng sai của các mệnh đề sau:}{%
  \y{Quãng đường vật đi được trong $4$ giây đầu là $164$ m.}{s}
  \y{Quãng đường vật đi được sau $t$ giây là $s(t)=12t+4$ (m).}{s}
  \y{Gia tốc của vật tại thời điểm $t=1$ giây là $10$ m/s$^2$.}{s}
  \y{Vận tốc trung bình của vật trong $4$ giây đầu là $40$ m/s.}{d}}

\cauDS[id=C4-B11-D04-TH-057,muc=TH,loigiai={$s(t)=\displaystyle\int v(t)\,\mathrm{d}t=2t^{3}+2t^{2}$ (vì $s(0)=0$); $a(t)=v'(t)=12t+4$. a) Sai: giá trị đúng là $72$ m. b) Đúng. c) Đúng. d) Đúng.}]{Một vật chuyển động thẳng với vận tốc $v(t)=6t^{2}+4t$ (m/s), $t$ là thời gian tính bằng giây, quãng đường tính từ thời điểm $t=0$. Xét tính đúng sai của các mệnh đề sau:}{%
  \y{Quãng đường vật đi được trong $3$ giây đầu là $75$ m.}{s}
  \y{Quãng đường vật đi được sau $t$ giây là $s(t)=2t^{3}+2t^{2}$ (m).}{d}
  \y{Gia tốc của vật tại thời điểm $t=1$ giây là $16$ m/s$^2$.}{d}
  \y{Vận tốc trung bình của vật trong $3$ giây đầu là $24$ m/s.}{d}}

\cauDS[id=C4-B11-D04-TH-058,muc=TH,loigiai={$s(t)=\displaystyle\int v(t)\,\mathrm{d}t=t^{3}+2t^{2}$ (vì $s(0)=0$); $a(t)=v'(t)=6t+4$. a) Đúng. b) Đúng. c) Sai: giá trị đúng là $96$ m. d) Sai: đó là vận tốc tại $t=1$.}]{Một vật chuyển động thẳng với vận tốc $v(t)=3t^{2}+4t$ (m/s), $t$ là thời gian tính bằng giây, quãng đường tính từ thời điểm $t=0$. Xét tính đúng sai của các mệnh đề sau:}{%
  \y{Quãng đường vật đi được trong $4$ giây đầu là $96$ m.}{d}
  \y{Vận tốc trung bình của vật trong $4$ giây đầu là $24$ m/s.}{d}
  \y{Quãng đường vật đi được trong $4$ giây đầu là $100$ m.}{s}
  \y{Gia tốc của vật tại thời điểm $t=1$ giây là $7$ m/s$^2$.}{s}}

\cauDS[id=C4-B11-D04-TH-059,muc=TH,loigiai={$s(t)=\displaystyle\int v(t)\,\mathrm{d}t=t^{3}+2t^{2}$ (vì $s(0)=0$); $a(t)=v'(t)=6t+4$. a) Đúng. b) Sai: đó là gia tốc. c) Đúng. d) Sai: đó là vận tốc tại $t=1$.}]{Một vật chuyển động thẳng với vận tốc $v(t)=3t^{2}+4t$ (m/s), $t$ là thời gian tính bằng giây, quãng đường tính từ thời điểm $t=0$. Xét tính đúng sai của các mệnh đề sau:}{%
  \y{Gia tốc của vật tại thời điểm $t=1$ giây là $10$ m/s$^2$.}{d}
  \y{Quãng đường vật đi được sau $t$ giây là $s(t)=6t+4$ (m).}{s}
  \y{Quãng đường vật đi được trong $3$ giây đầu là $45$ m.}{d}
  \y{Gia tốc của vật tại thời điểm $t=1$ giây là $7$ m/s$^2$.}{s}}

\cauDS[id=C4-B11-D04-TH-060,muc=TH,loigiai={$s(t)=\displaystyle\int v(t)\,\mathrm{d}t=t^{3}+t^{2}$ (vì $s(0)=0$); $a(t)=v'(t)=6t+2$. a) Đúng. b) Sai: đó là vận tốc tại $t=1$. c) Đúng. d) Sai: giá trị đúng là $80$ m.}]{Một vật chuyển động thẳng với vận tốc $v(t)=3t^{2}+2t$ (m/s), $t$ là thời gian tính bằng giây, quãng đường tính từ thời điểm $t=0$. Xét tính đúng sai của các mệnh đề sau:}{%
  \y{Quãng đường vật đi được sau $t$ giây là $s(t)=t^{3}+t^{2}$ (m).}{d}
  \y{Gia tốc của vật tại thời điểm $t=1$ giây là $5$ m/s$^2$.}{s}
  \y{Quãng đường vật đi được trong $4$ giây đầu là $80$ m.}{d}
  \y{Quãng đường vật đi được trong $4$ giây đầu là $84$ m.}{s}}

\cauDS[id=C4-B11-D04-TH-061,muc=TH,loigiai={$s(t)=\displaystyle\int v(t)\,\mathrm{d}t=2t^{3}+2t^{2}$ (vì $s(0)=0$); $a(t)=v'(t)=12t+4$. a) Sai: giá trị đúng là $160$ m. b) Đúng. c) Đúng. d) Sai: đó là vận tốc tại $t=1$.}]{Một vật chuyển động thẳng với vận tốc $v(t)=6t^{2}+4t$ (m/s), $t$ là thời gian tính bằng giây, quãng đường tính từ thời điểm $t=0$. Xét tính đúng sai của các mệnh đề sau:}{%
  \y{Quãng đường vật đi được trong $4$ giây đầu là $164$ m.}{s}
  \y{Gia tốc của vật tại thời điểm $t=1$ giây là $16$ m/s$^2$.}{d}
  \y{Vận tốc trung bình của vật trong $4$ giây đầu là $40$ m/s.}{d}
  \y{Gia tốc của vật tại thời điểm $t=1$ giây là $10$ m/s$^2$.}{s}}

\cauDS[id=C4-B11-D04-TH-062,muc=TH,loigiai={$s(t)=\displaystyle\int v(t)\,\mathrm{d}t=2t^{3}+2t^{2}$ (vì $s(0)=0$); $a(t)=v'(t)=12t+4$. a) Đúng. b) Đúng. c) Đúng. d) Sai: đó là vận tốc tại $t=1$.}]{Một vật chuyển động thẳng với vận tốc $v(t)=6t^{2}+4t$ (m/s), $t$ là thời gian tính bằng giây, quãng đường tính từ thời điểm $t=0$. Xét tính đúng sai của các mệnh đề sau:}{%
  \y{Quãng đường vật đi được trong $4$ giây đầu là $160$ m.}{d}
  \y{Vận tốc trung bình của vật trong $4$ giây đầu là $40$ m/s.}{d}
  \y{Gia tốc của vật tại thời điểm $t=1$ giây là $16$ m/s$^2$.}{d}
  \y{Gia tốc của vật tại thời điểm $t=1$ giây là $10$ m/s$^2$.}{s}}

\cauDS[id=C4-B11-D04-TH-063,muc=TH,loigiai={$s(t)=\displaystyle\int v(t)\,\mathrm{d}t=t^{3}+t^{2}$ (vì $s(0)=0$); $a(t)=v'(t)=6t+2$. a) Đúng. b) Sai: đó là vận tốc tại $t=1$. c) Sai: giá trị đúng là $12$ m. d) Đúng.}]{Một vật chuyển động thẳng với vận tốc $v(t)=3t^{2}+2t$ (m/s), $t$ là thời gian tính bằng giây, quãng đường tính từ thời điểm $t=0$. Xét tính đúng sai của các mệnh đề sau:}{%
  \y{Quãng đường vật đi được trong $2$ giây đầu là $12$ m.}{d}
  \y{Gia tốc của vật tại thời điểm $t=1$ giây là $5$ m/s$^2$.}{s}
  \y{Quãng đường vật đi được trong $2$ giây đầu là $14$ m.}{s}
  \y{Gia tốc của vật tại thời điểm $t=1$ giây là $8$ m/s$^2$.}{d}}

\cauDS[id=C4-B11-D04-TH-064,muc=TH,loigiai={$s(t)=\displaystyle\int v(t)\,\mathrm{d}t=t^{3}+t^{2}$ (vì $s(0)=0$); $a(t)=v'(t)=6t+2$. a) Sai: đó là gia tốc. b) Đúng. c) Đúng. d) Đúng.}]{Một vật chuyển động thẳng với vận tốc $v(t)=3t^{2}+2t$ (m/s), $t$ là thời gian tính bằng giây, quãng đường tính từ thời điểm $t=0$. Xét tính đúng sai của các mệnh đề sau:}{%
  \y{Quãng đường vật đi được sau $t$ giây là $s(t)=6t+2$ (m).}{s}
  \y{Quãng đường vật đi được sau $t$ giây là $s(t)=t^{3}+t^{2}$ (m).}{d}
  \y{Quãng đường vật đi được trong $4$ giây đầu là $80$ m.}{d}
  \y{Gia tốc của vật tại thời điểm $t=1$ giây là $8$ m/s$^2$.}{d}}

\cauTLN[id=C4-B11-D04-TH-065,muc=TH,loigiai={Gọi $S(t)$ là quãng đường, $S'(t)=v(t)$ nên $S(t)=t^{2} + 6 t+C$; $S(0)=0\Rightarrow C=0$. $S(25)=775$ (m).}]{Một máy bay chạy đà trên đường băng với vận tốc $v(t)=6+2t$ (m/s), trong đó $t$ là thời gian (giây) kể từ lúc bắt đầu chạy đà. Sau $25$ giây máy bay rời đường băng. Quãng đường (mét) máy bay đã chạy trên đường băng. Tính kết quả (đơn vị: m).}{775}

\cauTLN[id=C4-B11-D04-TH-066,muc=TH,loigiai={Lượng nước $V(t)=\displaystyle\int r(t)\,\mathrm{d}t=2t^2+30t$ (vì $V(0)=0$); $V(15)=900$ (lít).}]{Nước được bơm vào một bể với tốc độ $r(t)=4t+30$ (lít/phút), $t$ tính bằng phút. Lúc đầu bể không có nước. Sau $15$ phút, lượng nước trong bể. Tính kết quả (đơn vị: lít).}{900}

\cauTLN[id=C4-B11-D04-TH-067,muc=TH,loigiai={$v(t)=\displaystyle\int a(t)\,\mathrm{d}t=3t^2+t+C$, $v(0)=0\Rightarrow C=0$. $v(3)=30$ (m/s).}]{Một vật chuyển động với gia tốc $a(t)=6t+1$ (m/s$^2$), $t$ tính bằng giây. Biết vận tốc ban đầu của vật là $0$ m/s. Vận tốc của vật tại thời điểm $t=3$ giây. Tính kết quả (đơn vị: m/s).}{30}

\cauTLN[id=C4-B11-D04-TH-068,muc=TH,loigiai={Gọi $S(t)$ là quãng đường, $S'(t)=v(t)$ nên $S(t)=2 t^{2} + 8 t+C$; $S(0)=0\Rightarrow C=0$. $S(20)=960$ (m).}]{Một máy bay chạy đà trên đường băng với vận tốc $v(t)=8+4t$ (m/s), trong đó $t$ là thời gian (giây) kể từ lúc bắt đầu chạy đà. Sau $20$ giây máy bay rời đường băng. Quãng đường (mét) máy bay đã chạy trên đường băng. Tính kết quả (đơn vị: m).}{960}

\cauTLN[id=C4-B11-D04-TH-069,muc=TH,loigiai={$C(x)=1200\,000+50x-\dfrac{1}{10}x^2$ (nghìn đồng). $C(300)=1206000$ nghìn đồng $=1206$ triệu đồng.}]{Chi phí biên (nghìn đồng/sản phẩm) để sản xuất $x$ sản phẩm là $C'(x)=50-\dfrac{1}{5}x$. Chi phí cố định (khi $x=0$) là $1200$ triệu đồng. Tổng chi phí (triệu đồng) để sản xuất $300$ sản phẩm. Tính kết quả (đơn vị: triệu đồng).}{1206}

\cauTLN[id=C4-B11-D04-TH-070,muc=TH,loigiai={$P(t)=1000+40t+t^2$ (vì $P(0)=1000$). $P(10)=1500$ (con).}]{Một quần thể vi khuẩn ban đầu có $1000$ con. Tốc độ tăng trưởng của quần thể tại thời điểm $t$ (giờ) là $P'(t)=40+2t$ (con/giờ). Số vi khuẩn sau $10$ giờ. Tính kết quả (đơn vị: con).}{1500}

\cauTLN[id=C4-B11-D04-TH-071,muc=TH,loigiai={Lượng nước $V(t)=\displaystyle\int r(t)\,\mathrm{d}t=2t^2+30t$ (vì $V(0)=0$); $V(5)=200$ (lít).}]{Nước được bơm vào một bể với tốc độ $r(t)=4t+30$ (lít/phút), $t$ tính bằng phút. Lúc đầu bể không có nước. Sau $5$ phút, lượng nước trong bể. Tính kết quả (đơn vị: lít).}{200}

\cauTLN[id=C4-B11-D04-TH-072,muc=TH,loigiai={Xe dừng khi $v(t)=0\Leftrightarrow t=5$. $S(t)=\displaystyle\int v(t)\,\mathrm{d}t=-\dfrac{3}{2}t^2+15t$ (vì $S(0)=0$), $S(5)=\dfrac{75}{2}$ (m).}]{Một ô tô đang chạy với vận tốc $15$ m/s thì người lái đạp phanh; từ thời điểm đó ô tô chuyển động chậm dần đều với vận tốc $v(t)=-3t+15$ (m/s), $t$ là thời gian (giây) tính từ lúc đạp phanh. Từ lúc đạp phanh đến khi dừng hẳn, ô tô di chuyển được bao nhiêu mét?}{37,5}

\cauTLN[id=C4-B11-D04-TH-073,muc=TH,loigiai={$P(t)=1000+50t+t^2$ (vì $P(0)=1000$). $P(12)=1744$ (con).}]{Một quần thể vi khuẩn ban đầu có $1000$ con. Tốc độ tăng trưởng của quần thể tại thời điểm $t$ (giờ) là $P'(t)=50+2t$ (con/giờ). Số vi khuẩn sau $12$ giờ. Tính kết quả (đơn vị: con).}{1744}

\cauTLN[id=C4-B11-D04-TH-074,muc=TH,loigiai={Gọi $S(t)$ là quãng đường, $S'(t)=v(t)$ nên $S(t)=\dfrac{3 t^{2}}{2} + 8 t+C$; $S(0)=0\Rightarrow C=0$. $S(25)=\dfrac{2275}{2}$ (m).}]{Một máy bay chạy đà trên đường băng với vận tốc $v(t)=8+3t$ (m/s), trong đó $t$ là thời gian (giây) kể từ lúc bắt đầu chạy đà. Sau $25$ giây máy bay rời đường băng. Quãng đường (mét) máy bay đã chạy trên đường băng. Tính kết quả (đơn vị: m) (làm tròn kết quả đến hàng đơn vị).}{1138}

\cauTLN[id=C4-B11-D04-TH-075,muc=TH,loigiai={Xe dừng khi $v(t)=0\Leftrightarrow t=5$. $S(t)=\displaystyle\int v(t)\,\mathrm{d}t=-t^2+10t$ (vì $S(0)=0$), $S(5)=25$ (m).}]{Một ô tô đang chạy với vận tốc $10$ m/s thì người lái đạp phanh; từ thời điểm đó ô tô chuyển động chậm dần đều với vận tốc $v(t)=-2t+10$ (m/s), $t$ là thời gian (giây) tính từ lúc đạp phanh. Từ lúc đạp phanh đến khi dừng hẳn, ô tô di chuyển được bao nhiêu mét?}{25}

\cauTLN[id=C4-B11-D04-TH-076,muc=TH,loigiai={Gọi $S(t)$ là quãng đường, $S'(t)=v(t)$ nên $S(t)=\dfrac{3 t^{2}}{2} + 8 t+C$; $S(0)=0\Rightarrow C=0$. $S(20)=760$ (m).}]{Một máy bay chạy đà trên đường băng với vận tốc $v(t)=8+3t$ (m/s), trong đó $t$ là thời gian (giây) kể từ lúc bắt đầu chạy đà. Sau $20$ giây máy bay rời đường băng. Quãng đường (mét) máy bay đã chạy trên đường băng. Tính kết quả (đơn vị: m).}{760}

\cauTLN[id=C4-B11-D04-TH-077,muc=TH,loigiai={Lượng nước $V(t)=\displaystyle\int r(t)\,\mathrm{d}t=t^2+20t$ (vì $V(0)=0$); $V(10)=300$ (lít).}]{Nước được bơm vào một bể với tốc độ $r(t)=2t+20$ (lít/phút), $t$ tính bằng phút. Lúc đầu bể không có nước. Sau $10$ phút, lượng nước trong bể. Tính kết quả (đơn vị: lít).}{300}

\cauTLN[id=C4-B11-D04-TH-078,muc=TH,loigiai={Lượng nước $V(t)=\displaystyle\int r(t)\,\mathrm{d}t=3t^2+20t$ (vì $V(0)=0$); $V(10)=500$ (lít).}]{Nước được bơm vào một bể với tốc độ $r(t)=6t+20$ (lít/phút), $t$ tính bằng phút. Lúc đầu bể không có nước. Sau $10$ phút, lượng nước trong bể. Tính kết quả (đơn vị: lít).}{500}

\cauTLN[id=C4-B11-D04-TH-079,muc=TH,loigiai={Xe dừng khi $v(t)=0\Leftrightarrow t=3$. $S(t)=\displaystyle\int v(t)\,\mathrm{d}t=-3t^2+18t$ (vì $S(0)=0$), $S(3)=27$ (m).}]{Một ô tô đang chạy với vận tốc $18$ m/s thì người lái đạp phanh; từ thời điểm đó ô tô chuyển động chậm dần đều với vận tốc $v(t)=-6t+18$ (m/s), $t$ là thời gian (giây) tính từ lúc đạp phanh. Từ lúc đạp phanh đến khi dừng hẳn, ô tô di chuyển được bao nhiêu mét?}{27}

\cauTLN[id=C4-B11-D04-TH-080,muc=TH,loigiai={Lượng nước $V(t)=\displaystyle\int r(t)\,\mathrm{d}t=3t^2+10t$ (vì $V(0)=0$); $V(15)=825$ (lít).}]{Nước được bơm vào một bể với tốc độ $r(t)=6t+10$ (lít/phút), $t$ tính bằng phút. Lúc đầu bể không có nước. Sau $15$ phút, lượng nước trong bể. Tính kết quả (đơn vị: lít).}{825}

\cauTN[id=C4-B11-D04-VD-001,muc=VD,loigiai={$f(x)=x^{2}-\dfrac{2}{x}+C$; $f(1)=2\Rightarrow C=3$. $f(2)=6$.}]{Cho hàm số $y=f(x)$ xác định trên $(0;+\infty)$ có $f'(x)=2x+\dfrac{2}{x^2}$ với mọi $x>0$ và $f(1)=2$. Giá trị $f(2)$ bằng}
  {$6$}
  {$3$}{$10$}{$4$}

\cauTN[id=C4-B11-D04-VD-002,muc=VD,loigiai={$f(x)=x^{2}+\dfrac{1}{x}+C$; $f(1)=-2\Rightarrow C=-4$. $f(2)=\dfrac{1}{2}$.}]{Cho hàm số $y=f(x)$ xác định trên $(0;+\infty)$ có $f'(x)=2x-\dfrac{1}{x^2}$ với mọi $x>0$ và $f(1)=-2$. Giá trị $f(2)$ bằng}
  {$\dfrac{1}{2}$}
  {$\dfrac{9}{2}$}{$-\dfrac{3}{2}$}{$\dfrac{5}{2}$}

\cauTN[id=C4-B11-D04-VD-003,muc=VD,loigiai={$h(t)=49t-4{,}9t^2$; $v(t)=0\Leftrightarrow t=5$; $h_{\max}=h(5)=\dfrac{245}{2}\approx 123$ (m).}]{Một viên đạn được bắn thẳng đứng lên trên từ mặt đất với vận tốc $v(t)=49-9{,}8t$ (m/s). Độ cao lớn nhất (mét) mà viên đạn đạt được (làm tròn đến hàng đơn vị) là}
  {$123$ m}
  {$245$ m}{$61$ m}{$172$ m}

\cauTN[id=C4-B11-D04-VD-004,muc=VD,loigiai={$f(x)=3x^{2}+\dfrac{3}{x}+C$; $f(1)=2\Rightarrow C=-4$. $f(2)=\dfrac{19}{2}$.}]{Cho hàm số $y=f(x)$ xác định trên $(0;+\infty)$ có $f'(x)=6x-\dfrac{3}{x^2}$ với mọi $x>0$ và $f(1)=2$. Giá trị $f(2)$ bằng}
  {$\dfrac{19}{2}$}
  {$\dfrac{27}{2}$}{$\dfrac{7}{2}$}{$\dfrac{15}{2}$}

\cauTN[id=C4-B11-D04-VD-005,muc=VD,loigiai={Cá dừng khi $v(t)=0\Leftrightarrow t=15$. $s(t)=3t-\dfrac{1}{10}t^2$, $s(15)=\dfrac{45}{2}$ (km).}]{Cá hồi bơi ngược dòng sông với vận tốc $v(t)=-\dfrac{t}{5}+3$ (km/h), coi $t=0$ là lúc cá bắt đầu bơi vào dòng sông. Khoảng cách xa nhất (km) mà con cá có thể bơi được là}
  {$\dfrac{45}{2}$ km}
  {$45$ km}{$15$ km}{$\dfrac{45}{4}$ km}

\cauTN[id=C4-B11-D04-VD-006,muc=VD,loigiai={$h(t)=196t-4{,}9t^2$; $v(t)=0\Leftrightarrow t=20$; $h_{\max}=h(20)=1960\approx 1960$ (m).}]{Một viên đạn được bắn thẳng đứng lên trên từ mặt đất với vận tốc $v(t)=196-9{,}8t$ (m/s). Độ cao lớn nhất (mét) mà viên đạn đạt được (làm tròn đến hàng đơn vị) là}
  {$1960$ m}
  {$3920$ m}{$980$ m}{$2009$ m}

\cauTN[id=C4-B11-D04-VD-007,muc=VD,loigiai={$f(x)=3x^{2}-\dfrac{1}{x}+C$; $f(1)=-3\Rightarrow C=-5$. $f(2)=\dfrac{13}{2}$.}]{Cho hàm số $y=f(x)$ xác định trên $(0;+\infty)$ có $f'(x)=6x+\dfrac{1}{x^2}$ với mọi $x>0$ và $f(1)=-3$. Giá trị $f(2)$ bằng}
  {$\dfrac{13}{2}$}
  {$\dfrac{23}{2}$}{$\dfrac{17}{2}$}{$\dfrac{19}{2}$}

\cauTN[id=C4-B11-D04-VD-008,muc=VD,loigiai={Cá dừng khi $v(t)=0\Leftrightarrow t=15$. $s(t)=6t-\dfrac{2}{10}t^2$, $s(15)=45$ (km).}]{Cá hồi bơi ngược dòng sông với vận tốc $v(t)=-\dfrac{2t}{5}+6$ (km/h), coi $t=0$ là lúc cá bắt đầu bơi vào dòng sông. Khoảng cách xa nhất (km) mà con cá có thể bơi được là}
  {$45$ km}
  {$90$ km}{$15$ km}{$\dfrac{45}{2}$ km}

\cauTN[id=C4-B11-D04-VD-009,muc=VD,loigiai={$f(x)=x^{2}+\dfrac{2}{x}+C$; $f(1)=1\Rightarrow C=-2$. $f(4)=\dfrac{29}{2}$.}]{Cho hàm số $y=f(x)$ xác định trên $(0;+\infty)$ có $f'(x)=2x-\dfrac{2}{x^2}$ với mọi $x>0$ và $f(1)=1$. Giá trị $f(4)$ bằng}
  {$\dfrac{29}{2}$}
  {$\dfrac{33}{2}$}{$\dfrac{21}{2}$}{$\dfrac{27}{2}$}

\cauTN[id=C4-B11-D04-VD-010,muc=VD,loigiai={$f(x)=x^{2}+\dfrac{3}{x}+C$; $f(1)=-1\Rightarrow C=-5$. $f(2)=\dfrac{1}{2}$.}]{Cho hàm số $y=f(x)$ xác định trên $(0;+\infty)$ có $f'(x)=2x-\dfrac{3}{x^2}$ với mọi $x>0$ và $f(1)=-1$. Giá trị $f(2)$ bằng}
  {$\dfrac{1}{2}$}
  {$\dfrac{11}{2}$}{$-\dfrac{11}{2}$}{$\dfrac{3}{2}$}

\cauTN[id=C4-B11-D04-VD-011,muc=VD,loigiai={Cá dừng khi $v(t)=0\Leftrightarrow t=10$. $s(t)=8t-\dfrac{4}{10}t^2$, $s(10)=40$ (km).}]{Cá hồi bơi ngược dòng sông với vận tốc $v(t)=-\dfrac{4t}{5}+8$ (km/h), coi $t=0$ là lúc cá bắt đầu bơi vào dòng sông. Khoảng cách xa nhất (km) mà con cá có thể bơi được là}
  {$40$ km}
  {$80$ km}{$10$ km}{$20$ km}

\cauTN[id=C4-B11-D04-VD-012,muc=VD,loigiai={$f(x)=x^{2}+\dfrac{3}{x}+C$; $f(1)=-2\Rightarrow C=-6$. $f(4)=\dfrac{43}{4}$.}]{Cho hàm số $y=f(x)$ xác định trên $(0;+\infty)$ có $f'(x)=2x-\dfrac{3}{x^2}$ với mọi $x>0$ và $f(1)=-2$. Giá trị $f(4)$ bằng}
  {$\dfrac{43}{4}$}
  {$\dfrac{67}{4}$}{$\dfrac{19}{4}$}{$\dfrac{51}{4}$}

\cauTN[id=C4-B11-D04-VD-013,muc=VD,loigiai={$h(t)=160t-4{,}9t^2$; $v(t)=0\Leftrightarrow t=\dfrac{800}{49}$; $h_{\max}=h(\dfrac{800}{49})=\dfrac{64000}{49}\approx 1306$ (m).}]{Một viên đạn được bắn thẳng đứng lên trên từ mặt đất với vận tốc $v(t)=160-9{,}8t$ (m/s). Độ cao lớn nhất (mét) mà viên đạn đạt được (làm tròn đến hàng đơn vị) là}
  {$1306$ m}
  {$2612$ m}{$653$ m}{$1355$ m}

\cauTN[id=C4-B11-D04-VD-014,muc=VD,loigiai={$f(x)=2x^{2}+\dfrac{2}{x}+C$; $f(1)=2\Rightarrow C=-2$. $f(2)=7$.}]{Cho hàm số $y=f(x)$ xác định trên $(0;+\infty)$ có $f'(x)=4x-\dfrac{2}{x^2}$ với mọi $x>0$ và $f(1)=2$. Giá trị $f(2)$ bằng}
  {$7$}
  {$9$}{$3$}{$5$}

\cauTN[id=C4-B11-D04-VD-015,muc=VD,loigiai={$h(t)=98t-4{,}9t^2$; $v(t)=0\Leftrightarrow t=10$; $h_{\max}=h(10)=490\approx 490$ (m).}]{Một viên đạn được bắn thẳng đứng lên trên từ mặt đất với vận tốc $v(t)=98-9{,}8t$ (m/s). Độ cao lớn nhất (mét) mà viên đạn đạt được (làm tròn đến hàng đơn vị) là}
  {$490$ m}
  {$980$ m}{$245$ m}{$539$ m}

\cauTN[id=C4-B11-D04-VD-016,muc=VD,loigiai={$f(x)=3x^{2}-\dfrac{1}{x}+C$; $f(1)=2\Rightarrow C=0$. $f(2)=\dfrac{23}{2}$.}]{Cho hàm số $y=f(x)$ xác định trên $(0;+\infty)$ có $f'(x)=6x+\dfrac{1}{x^2}$ với mọi $x>0$ và $f(1)=2$. Giá trị $f(2)$ bằng}
  {$\dfrac{23}{2}$}
  {$\dfrac{27}{2}$}{$\dfrac{19}{2}$}{$12.5$}

\cauTN[id=C4-B11-D04-VD-017,muc=VD,loigiai={$h(t)=120t-4{,}9t^2$; $v(t)=0\Leftrightarrow t=\dfrac{600}{49}$; $h_{\max}=h(\dfrac{600}{49})=\dfrac{36000}{49}\approx 735$ (m).}]{Một viên đạn được bắn thẳng đứng lên trên từ mặt đất với vận tốc $v(t)=120-9{,}8t$ (m/s). Độ cao lớn nhất (mét) mà viên đạn đạt được (làm tròn đến hàng đơn vị) là}
  {$735$ m}
  {$1469$ m}{$367$ m}{$784$ m}

\cauTN[id=C4-B11-D04-VD-018,muc=VD,loigiai={$f(x)=x^{2}-\dfrac{3}{x}+C$; $f(1)=-3\Rightarrow C=-1$. $f(4)=\dfrac{57}{4}$.}]{Cho hàm số $y=f(x)$ xác định trên $(0;+\infty)$ có $f'(x)=2x+\dfrac{3}{x^2}$ với mọi $x>0$ và $f(1)=-3$. Giá trị $f(4)$ bằng}
  {$\dfrac{57}{4}$}
  {$\dfrac{61}{4}$}{$\dfrac{81}{4}$}{$\dfrac{69}{4}$}

\cauTN[id=C4-B11-D04-VD-019,muc=VD,loigiai={$f(x)=3x^{2}-\dfrac{1}{x}+C$; $f(1)=3\Rightarrow C=1$. $f(4)=\dfrac{195}{4}$.}]{Cho hàm số $y=f(x)$ xác định trên $(0;+\infty)$ có $f'(x)=6x+\dfrac{1}{x^2}$ với mọi $x>0$ và $f(1)=3$. Giá trị $f(4)$ bằng}
  {$\dfrac{195}{4}$}
  {$\dfrac{191}{4}$}{$\dfrac{203}{4}$}{$\dfrac{183}{4}$}

\cauTN[id=C4-B11-D04-VD-020,muc=VD,loigiai={$f(x)=x^{2}+\dfrac{3}{x}+C$; $f(1)=4\Rightarrow C=0$. $f(4)=\dfrac{67}{4}$.}]{Cho hàm số $y=f(x)$ xác định trên $(0;+\infty)$ có $f'(x)=2x-\dfrac{3}{x^2}$ với mọi $x>0$ và $f(1)=4$. Giá trị $f(4)$ bằng}
  {$\dfrac{67}{4}$}
  {$\dfrac{43}{4}$}{$\dfrac{51}{4}$}{$15.25$}

\cauTN[id=C4-B11-D04-VD-021,muc=VD,loigiai={$f(x)=3x^{2}-\dfrac{3}{x}+C$; $f(1)=-2\Rightarrow C=-2$. $f(2)=\dfrac{17}{2}$.}]{Cho hàm số $y=f(x)$ xác định trên $(0;+\infty)$ có $f'(x)=6x+\dfrac{3}{x^2}$ với mọi $x>0$ và $f(1)=-2$. Giá trị $f(2)$ bằng}
  {$\dfrac{17}{2}$}
  {$\dfrac{21}{2}$}{$\dfrac{29}{2}$}{$11.5$}

\cauTN[id=C4-B11-D04-VD-022,muc=VD,loigiai={$f(x)=2x^{2}-\dfrac{1}{x}+C$; $f(1)=2\Rightarrow C=1$. $f(2)=\dfrac{17}{2}$.}]{Cho hàm số $y=f(x)$ xác định trên $(0;+\infty)$ có $f'(x)=4x+\dfrac{1}{x^2}$ với mọi $x>0$ và $f(1)=2$. Giá trị $f(2)$ bằng}
  {$\dfrac{17}{2}$}
  {$\dfrac{15}{2}$}{$\dfrac{21}{2}$}{$\dfrac{13}{2}$}

\cauTN[id=C4-B11-D04-VD-023,muc=VD,loigiai={$h(t)=147t-4{,}9t^2$; $v(t)=0\Leftrightarrow t=15$; $h_{\max}=h(15)=\dfrac{2205}{2}\approx 1103$ (m).}]{Một viên đạn được bắn thẳng đứng lên trên từ mặt đất với vận tốc $v(t)=147-9{,}8t$ (m/s). Độ cao lớn nhất (mét) mà viên đạn đạt được (làm tròn đến hàng đơn vị) là}
  {$1103$ m}
  {$2205$ m}{$551$ m}{$1152$ m}

\cauTN[id=C4-B11-D04-VD-024,muc=VD,loigiai={$f(x)=2x^{2}-\dfrac{2}{x}+C$; $f(1)=4\Rightarrow C=4$. $f(2)=11$.}]{Cho hàm số $y=f(x)$ xác định trên $(0;+\infty)$ có $f'(x)=4x+\dfrac{2}{x^2}$ với mọi $x>0$ và $f(1)=4$. Giá trị $f(2)$ bằng}
  {$11$}
  {$7$}{$15$}{$13.0$}

\cauTN[id=C4-B11-D04-VD-025,muc=VD,loigiai={$f(x)=x^{2}-\dfrac{3}{x}+C$; $f(1)=1\Rightarrow C=3$. $f(4)=\dfrac{73}{4}$.}]{Cho hàm số $y=f(x)$ xác định trên $(0;+\infty)$ có $f'(x)=2x+\dfrac{3}{x^2}$ với mọi $x>0$ và $f(1)=1$. Giá trị $f(4)$ bằng}
  {$\dfrac{73}{4}$}
  {$\dfrac{61}{4}$}{$\dfrac{97}{4}$}{$\dfrac{69}{4}$}

\cauTN[id=C4-B11-D04-VD-026,muc=VD,loigiai={$f(x)=3x^{2}+\dfrac{1}{x}+C$; $f(1)=3\Rightarrow C=-1$. $f(2)=\dfrac{23}{2}$.}]{Cho hàm số $y=f(x)$ xác định trên $(0;+\infty)$ có $f'(x)=6x-\dfrac{1}{x^2}$ với mọi $x>0$ và $f(1)=3$. Giá trị $f(2)$ bằng}
  {$\dfrac{23}{2}$}
  {$\dfrac{25}{2}$}{$\dfrac{19}{2}$}{$\dfrac{17}{2}$}

\cauTN[id=C4-B11-D04-VD-027,muc=VD,loigiai={$f(x)=x^{2}+\dfrac{2}{x}+C$; $f(1)=-2\Rightarrow C=-5$. $f(4)=\dfrac{23}{2}$.}]{Cho hàm số $y=f(x)$ xác định trên $(0;+\infty)$ có $f'(x)=2x-\dfrac{2}{x^2}$ với mọi $x>0$ và $f(1)=-2$. Giá trị $f(4)$ bằng}
  {$\dfrac{23}{2}$}
  {$\dfrac{33}{2}$}{$\dfrac{15}{2}$}{$\dfrac{27}{2}$}

\cauTN[id=C4-B11-D04-VD-028,muc=VD,loigiai={$f(x)=3x^{2}+\dfrac{3}{x}+C$; $f(1)=3\Rightarrow C=-3$. $f(2)=\dfrac{21}{2}$.}]{Cho hàm số $y=f(x)$ xác định trên $(0;+\infty)$ có $f'(x)=6x-\dfrac{3}{x^2}$ với mọi $x>0$ và $f(1)=3$. Giá trị $f(2)$ bằng}
  {$\dfrac{21}{2}$}
  {$\dfrac{27}{2}$}{$\dfrac{9}{2}$}{$\dfrac{15}{2}$}

\cauTN[id=C4-B11-D04-VD-029,muc=VD,loigiai={$f(x)=x^{2}-\dfrac{3}{x}+C$; $f(1)=-1\Rightarrow C=1$. $f(4)=\dfrac{65}{4}$.}]{Cho hàm số $y=f(x)$ xác định trên $(0;+\infty)$ có $f'(x)=2x+\dfrac{3}{x^2}$ với mọi $x>0$ và $f(1)=-1$. Giá trị $f(4)$ bằng}
  {$\dfrac{65}{4}$}
  {$\dfrac{61}{4}$}{$\dfrac{89}{4}$}{$\dfrac{69}{4}$}

\cauTN[id=C4-B11-D04-VD-030,muc=VD,loigiai={$f(x)=2x^{2}+\dfrac{3}{x}+C$; $f(1)=3\Rightarrow C=-2$. $f(2)=\dfrac{15}{2}$.}]{Cho hàm số $y=f(x)$ xác định trên $(0;+\infty)$ có $f'(x)=4x-\dfrac{3}{x^2}$ với mọi $x>0$ và $f(1)=3$. Giá trị $f(2)$ bằng}
  {$\dfrac{15}{2}$}
  {$\dfrac{19}{2}$}{$\dfrac{3}{2}$}{$\dfrac{9}{2}$}

\cauTN[id=C4-B11-D04-VD-031,muc=VD,loigiai={$f(x)=x^{2}+\dfrac{2}{x}+C$; $f(1)=2\Rightarrow C=-1$. $f(4)=\dfrac{31}{2}$.}]{Cho hàm số $y=f(x)$ xác định trên $(0;+\infty)$ có $f'(x)=2x-\dfrac{2}{x^2}$ với mọi $x>0$ và $f(1)=2$. Giá trị $f(4)$ bằng}
  {$\dfrac{31}{2}$}
  {$\dfrac{33}{2}$}{$\dfrac{23}{2}$}{$\dfrac{27}{2}$}

\cauTN[id=C4-B11-D04-VD-032,muc=VD,loigiai={$f(x)=x^{2}-\dfrac{1}{x}+C$; $f(1)=2\Rightarrow C=2$. $f(2)=\dfrac{11}{2}$.}]{Cho hàm số $y=f(x)$ xác định trên $(0;+\infty)$ có $f'(x)=2x+\dfrac{1}{x^2}$ với mọi $x>0$ và $f(1)=2$. Giá trị $f(2)$ bằng}
  {$\dfrac{11}{2}$}
  {$\dfrac{7}{2}$}{$\dfrac{15}{2}$}{$6.5$}

\cauTN[id=C4-B11-D04-VD-033,muc=VD,loigiai={$f(x)=2x^{2}+\dfrac{3}{x}+C$; $f(1)=-1\Rightarrow C=-6$. $f(2)=\dfrac{7}{2}$.}]{Cho hàm số $y=f(x)$ xác định trên $(0;+\infty)$ có $f'(x)=4x-\dfrac{3}{x^2}$ với mọi $x>0$ và $f(1)=-1$. Giá trị $f(2)$ bằng}
  {$\dfrac{7}{2}$}
  {$\dfrac{19}{2}$}{$-\dfrac{5}{2}$}{$\dfrac{9}{2}$}

\cauTN[id=C4-B11-D04-VD-034,muc=VD,loigiai={$f(x)=3x^{2}+\dfrac{2}{x}+C$; $f(1)=-1\Rightarrow C=-6$. $f(2)=7$.}]{Cho hàm số $y=f(x)$ xác định trên $(0;+\infty)$ có $f'(x)=6x-\dfrac{2}{x^2}$ với mọi $x>0$ và $f(1)=-1$. Giá trị $f(2)$ bằng}
  {$7$}
  {$13$}{$3$}{$8$}

\cauTN[id=C4-B11-D04-VD-035,muc=VD,loigiai={$f(x)=3x^{2}-\dfrac{1}{x}+C$; $f(1)=-2\Rightarrow C=-4$. $f(2)=\dfrac{15}{2}$.}]{Cho hàm số $y=f(x)$ xác định trên $(0;+\infty)$ có $f'(x)=6x+\dfrac{1}{x^2}$ với mọi $x>0$ và $f(1)=-2$. Giá trị $f(2)$ bằng}
  {$\dfrac{15}{2}$}
  {$\dfrac{23}{2}$}{$\dfrac{19}{2}$}{$8.5$}

\cauTN[id=C4-B11-D04-VD-036,muc=VD,loigiai={$f(x)=x^{2}+\dfrac{2}{x}+C$; $f(1)=-3\Rightarrow C=-6$. $f(4)=\dfrac{21}{2}$.}]{Cho hàm số $y=f(x)$ xác định trên $(0;+\infty)$ có $f'(x)=2x-\dfrac{2}{x^2}$ với mọi $x>0$ và $f(1)=-3$. Giá trị $f(4)$ bằng}
  {$\dfrac{21}{2}$}
  {$\dfrac{33}{2}$}{$\dfrac{13}{2}$}{$\dfrac{27}{2}$}

\cauDS[id=C4-B11-D04-VD-037,muc=VD,loigiai={$h(t)=\displaystyle\int v(t)\,\mathrm{d}t=147t-4{,}9t^2$ (vì $h(0)=0$); $v(t)=0\Leftrightarrow t=15$. a) Sai: giá trị đúng là $1102{,}5$ m. b) Sai: thời gian đúng là $30$ giây. c) Sai: hệ số đúng là $4{,}9$. d) Đúng.}]{Một viên đạn được bắn thẳng đứng lên trên từ mặt đất với vận tốc $v(t)=147-9{,}8t$ (m/s), $t$ tính bằng giây ($t=0$ là thời điểm bắn). Xét tính đúng sai của các mệnh đề sau:}{%
  \y{Độ cao lớn nhất của viên đạn là $2205$ m.}{s}
  \y{Viên đạn rơi trở lại mặt đất sau $15$ giây kể từ lúc bắn.}{s}
  \y{Độ cao của viên đạn tại thời điểm $t$ là $h(t)=147t-9{,}8t^2$ (m).}{s}
  \y{Độ cao lớn nhất của viên đạn là $1102{,}5$ m.}{d}}

\cauDS[id=C4-B11-D04-VD-038,muc=VD,loigiai={$h(t)=\displaystyle\int v(t)\,\mathrm{d}t=147t-4{,}9t^2$ (vì $h(0)=0$); $v(t)=0\Leftrightarrow t=15$. a) Sai: thời gian đúng là $30$ giây. b) Đúng. c) Sai: giá trị đúng là $1102{,}5$ m. d) Sai: hệ số đúng là $4{,}9$.}]{Một viên đạn được bắn thẳng đứng lên trên từ mặt đất với vận tốc $v(t)=147-9{,}8t$ (m/s), $t$ tính bằng giây ($t=0$ là thời điểm bắn). Xét tính đúng sai của các mệnh đề sau:}{%
  \y{Viên đạn rơi trở lại mặt đất sau $15$ giây kể từ lúc bắn.}{s}
  \y{Độ cao của viên đạn tại thời điểm $t$ là $h(t)=147t-4{,}9t^2$ (m).}{d}
  \y{Độ cao lớn nhất của viên đạn là $2205$ m.}{s}
  \y{Độ cao của viên đạn tại thời điểm $t$ là $h(t)=147t-9{,}8t^2$ (m).}{s}}

\cauDS[id=C4-B11-D04-VD-039,muc=VD,loigiai={$h(t)=\displaystyle\int v(t)\,\mathrm{d}t=147t-4{,}9t^2$ (vì $h(0)=0$); $v(t)=0\Leftrightarrow t=15$. a) Đúng. b) Sai: giá trị đúng là $1102{,}5$ m. c) Đúng. d) Sai: hệ số đúng là $4{,}9$.}]{Một viên đạn được bắn thẳng đứng lên trên từ mặt đất với vận tốc $v(t)=147-9{,}8t$ (m/s), $t$ tính bằng giây ($t=0$ là thời điểm bắn). Xét tính đúng sai của các mệnh đề sau:}{%
  \y{Sau $3$ giây, viên đạn ở độ cao $396{,}9$ m.}{d}
  \y{Độ cao lớn nhất của viên đạn là $2205$ m.}{s}
  \y{Viên đạn đạt độ cao lớn nhất tại thời điểm $t=15$ giây.}{d}
  \y{Độ cao của viên đạn tại thời điểm $t$ là $h(t)=147t-9{,}8t^2$ (m).}{s}}

\cauDS[id=C4-B11-D04-VD-040,muc=VD,loigiai={$h(t)=\displaystyle\int v(t)\,\mathrm{d}t=294t-4{,}9t^2$ (vì $h(0)=0$); $v(t)=0\Leftrightarrow t=30$. a) Đúng. b) Sai: giá trị đúng là $4410$ m. c) Sai: hệ số đúng là $4{,}9$. d) Đúng.}]{Một viên đạn được bắn thẳng đứng lên trên từ mặt đất với vận tốc $v(t)=294-9{,}8t$ (m/s), $t$ tính bằng giây ($t=0$ là thời điểm bắn). Xét tính đúng sai của các mệnh đề sau:}{%
  \y{Sau $3$ giây, viên đạn ở độ cao $837{,}9$ m.}{d}
  \y{Độ cao lớn nhất của viên đạn là $8820$ m.}{s}
  \y{Độ cao của viên đạn tại thời điểm $t$ là $h(t)=294t-9{,}8t^2$ (m).}{s}
  \y{Độ cao của viên đạn tại thời điểm $t$ là $h(t)=294t-4{,}9t^2$ (m).}{d}}

\cauDS[id=C4-B11-D04-VD-041,muc=VD,loigiai={$h(t)=\displaystyle\int v(t)\,\mathrm{d}t=147t-4{,}9t^2$ (vì $h(0)=0$); $v(t)=0\Leftrightarrow t=15$. a) Đúng. b) Sai: hệ số đúng là $4{,}9$. c) Đúng. d) Sai: thời gian đúng là $30$ giây.}]{Một viên đạn được bắn thẳng đứng lên trên từ mặt đất với vận tốc $v(t)=147-9{,}8t$ (m/s), $t$ tính bằng giây ($t=0$ là thời điểm bắn). Xét tính đúng sai của các mệnh đề sau:}{%
  \y{Độ cao của viên đạn tại thời điểm $t$ là $h(t)=147t-4{,}9t^2$ (m).}{d}
  \y{Độ cao của viên đạn tại thời điểm $t$ là $h(t)=147t-9{,}8t^2$ (m).}{s}
  \y{Viên đạn đạt độ cao lớn nhất tại thời điểm $t=15$ giây.}{d}
  \y{Viên đạn rơi trở lại mặt đất sau $15$ giây kể từ lúc bắn.}{s}}

\cauDS[id=C4-B11-D04-VD-042,muc=VD,loigiai={$h(t)=\displaystyle\int v(t)\,\mathrm{d}t=147t-4{,}9t^2$ (vì $h(0)=0$); $v(t)=0\Leftrightarrow t=15$. a) Sai: thời gian đúng là $30$ giây. b) Sai: hệ số đúng là $4{,}9$. c) Sai: giá trị đúng là $1102{,}5$ m. d) Đúng.}]{Một viên đạn được bắn thẳng đứng lên trên từ mặt đất với vận tốc $v(t)=147-9{,}8t$ (m/s), $t$ tính bằng giây ($t=0$ là thời điểm bắn). Xét tính đúng sai của các mệnh đề sau:}{%
  \y{Viên đạn rơi trở lại mặt đất sau $15$ giây kể từ lúc bắn.}{s}
  \y{Độ cao của viên đạn tại thời điểm $t$ là $h(t)=147t-9{,}8t^2$ (m).}{s}
  \y{Độ cao lớn nhất của viên đạn là $2205$ m.}{s}
  \y{Sau $3$ giây, viên đạn ở độ cao $396{,}9$ m.}{d}}

\cauDS[id=C4-B11-D04-VD-043,muc=VD,loigiai={$h(t)=\displaystyle\int v(t)\,\mathrm{d}t=294t-4{,}9t^2$ (vì $h(0)=0$); $v(t)=0\Leftrightarrow t=30$. a) Đúng. b) Sai: giá trị đúng là $4410$ m. c) Đúng. d) Sai: hệ số đúng là $4{,}9$.}]{Một viên đạn được bắn thẳng đứng lên trên từ mặt đất với vận tốc $v(t)=294-9{,}8t$ (m/s), $t$ tính bằng giây ($t=0$ là thời điểm bắn). Xét tính đúng sai của các mệnh đề sau:}{%
  \y{Độ cao lớn nhất của viên đạn là $4410$ m.}{d}
  \y{Độ cao lớn nhất của viên đạn là $8820$ m.}{s}
  \y{Sau $1$ giây, viên đạn ở độ cao $289{,}1$ m.}{d}
  \y{Độ cao của viên đạn tại thời điểm $t$ là $h(t)=294t-9{,}8t^2$ (m).}{s}}

\cauDS[id=C4-B11-D04-VD-044,muc=VD,loigiai={$h(t)=\displaystyle\int v(t)\,\mathrm{d}t=294t-4{,}9t^2$ (vì $h(0)=0$); $v(t)=0\Leftrightarrow t=30$. a) Đúng. b) Sai: giá trị đúng là $4410$ m. c) Đúng. d) Sai: thời gian đúng là $60$ giây.}]{Một viên đạn được bắn thẳng đứng lên trên từ mặt đất với vận tốc $v(t)=294-9{,}8t$ (m/s), $t$ tính bằng giây ($t=0$ là thời điểm bắn). Xét tính đúng sai của các mệnh đề sau:}{%
  \y{Độ cao lớn nhất của viên đạn là $4410$ m.}{d}
  \y{Độ cao lớn nhất của viên đạn là $8820$ m.}{s}
  \y{Viên đạn đạt độ cao lớn nhất tại thời điểm $t=30$ giây.}{d}
  \y{Viên đạn rơi trở lại mặt đất sau $30$ giây kể từ lúc bắn.}{s}}

\cauDS[id=C4-B11-D04-VD-045,muc=VD,loigiai={$h(t)=\displaystyle\int v(t)\,\mathrm{d}t=49t-4{,}9t^2$ (vì $h(0)=0$); $v(t)=0\Leftrightarrow t=5$. a) Đúng. b) Sai: giá trị đúng là $122{,}5$ m. c) Đúng. d) Sai: hệ số đúng là $4{,}9$.}]{Một viên đạn được bắn thẳng đứng lên trên từ mặt đất với vận tốc $v(t)=49-9{,}8t$ (m/s), $t$ tính bằng giây ($t=0$ là thời điểm bắn). Xét tính đúng sai của các mệnh đề sau:}{%
  \y{Độ cao lớn nhất của viên đạn là $122{,}5$ m.}{d}
  \y{Độ cao lớn nhất của viên đạn là $245$ m.}{s}
  \y{Độ cao của viên đạn tại thời điểm $t$ là $h(t)=49t-4{,}9t^2$ (m).}{d}
  \y{Độ cao của viên đạn tại thời điểm $t$ là $h(t)=49t-9{,}8t^2$ (m).}{s}}

\cauDS[id=C4-B11-D04-VD-046,muc=VD,loigiai={$h(t)=\displaystyle\int v(t)\,\mathrm{d}t=196t-4{,}9t^2$ (vì $h(0)=0$); $v(t)=0\Leftrightarrow t=20$. a) Đúng. b) Sai: giá trị đúng là $1960$ m. c) Sai: hệ số đúng là $4{,}9$. d) Đúng.}]{Một viên đạn được bắn thẳng đứng lên trên từ mặt đất với vận tốc $v(t)=196-9{,}8t$ (m/s), $t$ tính bằng giây ($t=0$ là thời điểm bắn). Xét tính đúng sai của các mệnh đề sau:}{%
  \y{Độ cao lớn nhất của viên đạn là $1960$ m.}{d}
  \y{Độ cao lớn nhất của viên đạn là $3920$ m.}{s}
  \y{Độ cao của viên đạn tại thời điểm $t$ là $h(t)=196t-9{,}8t^2$ (m).}{s}
  \y{Độ cao của viên đạn tại thời điểm $t$ là $h(t)=196t-4{,}9t^2$ (m).}{d}}

\cauDS[id=C4-B11-D04-VD-047,muc=VD,loigiai={$h(t)=\displaystyle\int v(t)\,\mathrm{d}t=294t-4{,}9t^2$ (vì $h(0)=0$); $v(t)=0\Leftrightarrow t=30$. a) Đúng. b) Sai: thời gian đúng là $60$ giây. c) Đúng. d) Sai: giá trị đúng là $4410$ m.}]{Một viên đạn được bắn thẳng đứng lên trên từ mặt đất với vận tốc $v(t)=294-9{,}8t$ (m/s), $t$ tính bằng giây ($t=0$ là thời điểm bắn). Xét tính đúng sai của các mệnh đề sau:}{%
  \y{Viên đạn đạt độ cao lớn nhất tại thời điểm $t=30$ giây.}{d}
  \y{Viên đạn rơi trở lại mặt đất sau $30$ giây kể từ lúc bắn.}{s}
  \y{Sau $2$ giây, viên đạn ở độ cao $568{,}4$ m.}{d}
  \y{Độ cao lớn nhất của viên đạn là $8820$ m.}{s}}

\cauDS[id=C4-B11-D04-VD-048,muc=VD,loigiai={$h(t)=\displaystyle\int v(t)\,\mathrm{d}t=196t-4{,}9t^2$ (vì $h(0)=0$); $v(t)=0\Leftrightarrow t=20$. a) Sai: giá trị đúng là $1960$ m. b) Đúng. c) Đúng. d) Sai: hệ số đúng là $4{,}9$.}]{Một viên đạn được bắn thẳng đứng lên trên từ mặt đất với vận tốc $v(t)=196-9{,}8t$ (m/s), $t$ tính bằng giây ($t=0$ là thời điểm bắn). Xét tính đúng sai của các mệnh đề sau:}{%
  \y{Độ cao lớn nhất của viên đạn là $3920$ m.}{s}
  \y{Sau $4$ giây, viên đạn ở độ cao $705{,}6$ m.}{d}
  \y{Độ cao lớn nhất của viên đạn là $1960$ m.}{d}
  \y{Độ cao của viên đạn tại thời điểm $t$ là $h(t)=196t-9{,}8t^2$ (m).}{s}}

\cauTLN[id=C4-B11-D04-VD-049,muc=VD,loigiai={$f(x)=2x^{2}-\dfrac{3}{x}+C$, $C=0$, $f(2)=\dfrac{13}{2}$.}]{Cho hàm số $y=f(x)$ xác định trên $(0;+\infty)$ có $f'(x)=4x+\dfrac{3}{x^2}$ với mọi $x>0$ và $f(1)=-1$. Tính $f(2)$.}{6,5}

\cauTLN[id=C4-B11-D04-VD-050,muc=VD,loigiai={$h(t)=100t-4{,}9t^2$; độ cao lớn nhất tại $t=\dfrac{500}{49}$: $h=\dfrac{25000}{49}\approx 510$.}]{Một viên đạn được bắn thẳng đứng lên trên từ mặt đất với vận tốc $v(t)=100-9{,}8t$ (m/s). Tính độ cao lớn nhất (mét) mà viên đạn đạt được (làm tròn kết quả đến hàng đơn vị).}{510}

\cauTLN[id=C4-B11-D04-VD-051,muc=VD,loigiai={$h(t)=160t-4{,}9t^2$; độ cao lớn nhất tại $t=\dfrac{800}{49}$: $h=\dfrac{64000}{49}\approx 1306$.}]{Một viên đạn được bắn thẳng đứng lên trên từ mặt đất với vận tốc $v(t)=160-9{,}8t$ (m/s). Tính độ cao lớn nhất (mét) mà viên đạn đạt được (làm tròn kết quả đến hàng đơn vị).}{1306}

\cauTLN[id=C4-B11-D04-VD-052,muc=VD,loigiai={$f(x)=x^{2}+\dfrac{3}{x}+C$, $C=-4$, $f(4)=\dfrac{51}{4}$.}]{Cho hàm số $y=f(x)$ xác định trên $(0;+\infty)$ có $f'(x)=2x-\dfrac{3}{x^2}$ với mọi $x>0$ và $f(1)=0$. Tính $f(4)$ (làm tròn kết quả đến hàng phần mười).}{12,8}

\cauTLN[id=C4-B11-D04-VD-053,muc=VD,loigiai={$h(t)=150t-4{,}9t^2$; độ cao lớn nhất tại $t=\dfrac{750}{49}$: $h=\dfrac{56250}{49}\approx 1148$.}]{Một viên đạn được bắn thẳng đứng lên trên từ mặt đất với vận tốc $v(t)=150-9{,}8t$ (m/s). Tính độ cao lớn nhất (mét) mà viên đạn đạt được (làm tròn kết quả đến hàng phần mười).}{1148}

\cauTLN[id=C4-B11-D04-VD-054,muc=VD,loigiai={$f(x)=3x^{2}+\dfrac{2}{x}+C$, $C=-8$, $f(4)=\dfrac{81}{2}$.}]{Cho hàm số $y=f(x)$ xác định trên $(0;+\infty)$ có $f'(x)=6x-\dfrac{2}{x^2}$ với mọi $x>0$ và $f(1)=-3$. Tính $f(4)$.}{40,5}

\cauTLN[id=C4-B11-D04-VD-055,muc=VD,loigiai={$f(x)=2x^{2}+\dfrac{1}{x}+C$, $C=1$, $f(2)=\dfrac{19}{2}$.}]{Cho hàm số $y=f(x)$ xác định trên $(0;+\infty)$ có $f'(x)=4x-\dfrac{1}{x^2}$ với mọi $x>0$ và $f(1)=4$. Tính $f(2)$.}{9,5}

\cauTLN[id=C4-B11-D04-VD-056,muc=VD,loigiai={$f(x)=x^{2}-\dfrac{2}{x}+C$, $C=-1$, $f(4)=\dfrac{29}{2}$.}]{Cho hàm số $y=f(x)$ xác định trên $(0;+\infty)$ có $f'(x)=2x+\dfrac{2}{x^2}$ với mọi $x>0$ và $f(1)=-2$. Tính $f(4)$.}{14,5}

\cauTLN[id=C4-B11-D04-VD-057,muc=VD,loigiai={$f(x)=3x^{2}+\dfrac{3}{x}+C$, $C=-9$, $f(2)=\dfrac{9}{2}$.}]{Cho hàm số $y=f(x)$ xác định trên $(0;+\infty)$ có $f'(x)=6x-\dfrac{3}{x^2}$ với mọi $x>0$ và $f(1)=-3$. Tính $f(2)$.}{4,5}

\cauTLN[id=C4-B11-D04-VD-058,muc=VD,loigiai={$f(x)=2x^{2}-\dfrac{3}{x}+C$, $C=5$, $f(2)=\dfrac{23}{2}$.}]{Cho hàm số $y=f(x)$ xác định trên $(0;+\infty)$ có $f'(x)=4x+\dfrac{3}{x^2}$ với mọi $x>0$ và $f(1)=4$. Tính $f(2)$.}{11,5}

\cauTLN[id=C4-B11-D04-VD-059,muc=VD,loigiai={$f(x)=2x^{2}-\dfrac{1}{x}+C$, $C=1$, $f(4)=\dfrac{131}{4}$.}]{Cho hàm số $y=f(x)$ xác định trên $(0;+\infty)$ có $f'(x)=4x+\dfrac{1}{x^2}$ với mọi $x>0$ và $f(1)=2$. Tính $f(4)$ (làm tròn kết quả đến hàng phần mười).}{32,8}

\cauTLN[id=C4-B11-D04-VD-060,muc=VD,loigiai={$f(x)=3x^{2}-\dfrac{1}{x}+C$, $C=2$, $f(4)=\dfrac{199}{4}$.}]{Cho hàm số $y=f(x)$ xác định trên $(0;+\infty)$ có $f'(x)=6x+\dfrac{1}{x^2}$ với mọi $x>0$ và $f(1)=4$. Tính $f(4)$ (làm tròn kết quả đến hàng phần mười).}{49,8}
